\documentclass[12pt,a4paper,oneside]{report}

% ============================================================
% Rapport de projet ComptaFlow
% Les champs inconnus sont volontairement laissés vides.
% Compilation conseillée : pdflatex (deux passes) ou latexmk.
% ============================================================

\usepackage[utf8]{inputenc}
\usepackage[T1]{fontenc}
\usepackage[english,french]{babel}
\usepackage{lmodern}
\usepackage{microtype}
\usepackage[a4paper,margin=2.5cm,headheight=26pt]{geometry}
\usepackage{setspace}
\usepackage{graphicx}
\usepackage{float}
\usepackage{caption}
\usepackage{subcaption}
\usepackage{booktabs}
\usepackage{longtable}
\usepackage{tabularx}
\usepackage{array}
\usepackage{enumitem}
\usepackage[table]{xcolor}
\usepackage{fancyhdr}
\usepackage{titlesec}
\usepackage{hyperref}
\usepackage{url}
\usepackage{listings}
\usepackage{tikz}
\usetikzlibrary{arrows.meta,positioning,fit,shapes.geometric,shapes.multipart,calc}

\definecolor{ComptaGreen}{HTML}{1F5E45}
\definecolor{ComptaDark}{HTML}{111827}
\definecolor{ComptaLight}{HTML}{F7F9FB}
\definecolor{ComptaBlue}{HTML}{17324D}
\definecolor{ComptaSlate}{HTML}{52677D}
\definecolor{ComptaGray}{HTML}{64748B}
\definecolor{ComptaTableLine}{HTML}{718096}
\definecolor{ComptaTableRow}{HTML}{F4F6F8}

\hypersetup{
  colorlinks=true,
  linkcolor=ComptaGreen,
  citecolor=ComptaBlue,
  urlcolor=ComptaBlue,
  pdftitle={Conception et réalisation de ComptaFlow},
  pdfauthor={}
}

\setstretch{1.15}
\setlength{\parindent}{0pt}
\setlength{\parskip}{0.55em}
\setcounter{secnumdepth}{3}
\setcounter{tocdepth}{3}

\titleformat{\chapter}[display]
  {\normalfont\bfseries\color{ComptaDark}}
  {\Large\color{ComptaBlue}\chaptername\ \thechapter}
  {0.55em}{\Huge}
  [\vspace{0.35em}\color{ComptaBlue}\titlerule]
\titleformat{name=\chapter,numberless}[display]
  {\normalfont\bfseries\color{ComptaDark}}
  {}{0pt}{\Huge}
  [\vspace{0.35em}\color{ComptaBlue}\titlerule]
\titleformat{\section}
  {\normalfont\Large\bfseries\color{ComptaBlue}}
  {\thesection}{0.75em}{}
\titleformat{\subsection}
  {\normalfont\large\bfseries\color{ComptaGreen}}
  {\thesubsection}{0.75em}{}

\pagestyle{fancy}
\fancyhf{}
\fancyhead[L]{\textcolor{ComptaDark}{ComptaFlow}}
\fancyhead[R]{\textcolor{ComptaDark}{\nouppercase{\leftmark}}}
\fancyfoot[C]{\textcolor{ComptaGreen}{\thepage}}
\renewcommand{\chaptermark}[1]{\markboth{#1}{}}
\renewcommand{\headrulewidth}{0pt}
\renewcommand{\headrule}{\hbox to\headwidth{\color{ComptaBlue}\leaders\hrule height 0.5pt\hfill}}

\captionsetup{
  labelfont={color=ComptaBlue},
  textfont={color=ComptaDark}
}

\newcolumntype{Y}{>{\raggedright\arraybackslash}X}
\newcolumntype{L}[1]{>{\raggedright\arraybackslash}p{#1}}
\newcommand{\TableSetup}{%
  \arrayrulecolor{ComptaTableLine}%
  \setlength{\arrayrulewidth}{0.45pt}%
  \renewcommand{\arraystretch}{1.18}%
}
\newcommand{\TableHead}[1]{\textcolor{white}{\textbf{#1}}}
\newcommand{\TableId}[1]{\textcolor{ComptaGreen}{\textbf{#1}}}
\newcommand{\TableSuccess}[1]{\textcolor{ComptaGreen}{\textbf{#1}}}

% ---------- Informations du rapport ----------
\newcommand{\Auteur}{RADOUI Fatima Zahra}
\newcommand{\Formation}{Génie informatique}
\newcommand{\Etablissement}{ENSA de Berrechid}
\newcommand{\EntrepriseAccueil}{SEGURIBAT}
\newcommand{\EncadrantPedagogique}{M. Dehaj}
\newcommand{\EncadrantEntreprise}{M. Faiz Mohammed Er-Raitab}
\newcommand{\AnneeUniversitaire}{2025--2026}
\newcommand{\DateSoutenance}{}
\newcommand{\PeriodeProjet}{du 1\textsuperscript{er} au 30 juillet 2026}

% ---------- Textes modifiables de la page de garde ----------
\newcommand{\NomEcoleComplet}{ÉCOLE NATIONALE DES SCIENCES APPLIQUÉES DE BERRECHID}
\newcommand{\TypeRapport}{RAPPORT DE STAGE}
\newcommand{\SousTitreRapport}{Conception et réalisation d'une plateforme\\
  intelligente de gestion documentaire et de pré-comptabilité}
\newcommand{\MentionSeguribat}{COMPTABLE AGRÉÉ PAR L'ÉTAT}

% ---------- Logos et réglages modifiables de la page de garde ----------
% Les fichiers restent séparés : leur taille et leur position peuvent être
% adaptés ici sans transformer la couverture en image unique.
\newcommand{\LogoENSA}{assets/logos/ensa-berrechid-transparent.png}
\newcommand{\LogoUniversite}{assets/logos/universite-hassan-1er-transparent.png}
\newlength{\LogoENSAWidth}
\newlength{\LogoUniversiteHeight}
\newlength{\LogoENSAOffsetY}
\newlength{\LogoUniversiteOffsetY}
\setlength{\LogoENSAWidth}{6.7cm}
\setlength{\LogoUniversiteHeight}{3.05cm}
\setlength{\LogoENSAOffsetY}{0cm}
\setlength{\LogoUniversiteOffsetY}{-0.12cm}

% ---------- Espacements modifiables de la page de garde ----------
% Modifier ces valeurs permet de déplacer verticalement les grands blocs.
\newlength{\CouvertureEspaceApresLogos}
\newlength{\CouvertureEspaceApresEntreprise}
\newlength{\CouvertureEspaceAvantTitre}
\newlength{\CouvertureEspaceAvantInformations}
\setlength{\CouvertureEspaceApresLogos}{1.15cm}
\setlength{\CouvertureEspaceApresEntreprise}{1.25cm}
\setlength{\CouvertureEspaceAvantTitre}{0.75cm}
\setlength{\CouvertureEspaceAvantInformations}{0.8cm}

\newcommand{\capture}[3]{%
  \begin{figure}[H]
    \centering
    \IfFileExists{#1}{%
      \includegraphics[width=0.92\textwidth,height=0.62\textheight,keepaspectratio]{#1}%
    }{%
      \fcolorbox{ComptaGray}{white}{%
        \begin{minipage}[c][5.3cm][c]{0.86\textwidth}
          \centering
          \textbf{Illustration indisponible}\\[0.7em]
          Fichier attendu : \texttt{\detokenize{#1}}\\[0.5em]
          #3
        \end{minipage}%
      }%
    }
    \caption{#2}
  \end{figure}%
}

\tikzset{
  actor/.style={draw, rounded corners, fill=ComptaLight, minimum width=2.7cm,
    minimum height=0.9cm, align=center, font=\small\bfseries},
  usecase/.style={draw, ellipse, fill=white, minimum width=3.5cm,
    minimum height=1cm, align=center, font=\small},
  box/.style={draw, rounded corners, fill=white, minimum width=3.1cm,
    minimum height=1cm, align=center, font=\small},
  layer/.style={draw, rounded corners, very thick, minimum width=12.5cm,
    minimum height=1.15cm, align=center, font=\small},
  flow/.style={-{Latex[length=2.3mm]}, thick, draw=ComptaDark},
  data/.style={draw, rectangle split, rectangle split parts=2, rounded corners,
    text width=3.25cm, align=left, font=\scriptsize, fill=white},
  meriseEntity/.style={draw, rectangle split, rectangle split parts=2,
    rounded corners, text width=3.15cm, align=left, font=\scriptsize,
    fill=blue!4},
  meriseAssoc/.style={draw, ellipse, align=center, font=\scriptsize,
    fill=ComptaBlue!5, minimum width=2.2cm, minimum height=0.8cm},
  note/.style={draw, dashed, rounded corners, fill=ComptaGray!6, align=left, font=\scriptsize}
}

\lstset{
  basicstyle=\ttfamily\small,
  breaklines=true,
  frame=single,
  rulecolor=\color{ComptaGray},
  backgroundcolor=\color{black!2},
  showstringspaces=false
}

\begin{document}
\pagecolor{white}
\color{ComptaDark}

% ============================================================
% Page de garde
% ============================================================
\hypersetup{pageanchor=false}
\begin{titlepage}
  \newgeometry{left=2.1cm,right=2.1cm,top=1.25cm,bottom=1.15cm}
  \thispagestyle{empty}
  \sffamily

  % Logos supérieurs sur fond transparent. Les commandes du préambule
  % permettent de modifier séparément fichiers, dimensions et décalages.
  \noindent
  \begin{minipage}[t]{0.56\textwidth}
    \vspace{0pt}
    \raggedright
    \raisebox{\LogoENSAOffsetY}{%
      \includegraphics[width=\LogoENSAWidth,keepaspectratio]{\LogoENSA}}
  \end{minipage}%
  \hfill
  \begin{minipage}[t]{0.30\textwidth}
    \vspace{0pt}
    \raggedleft
    \raisebox{\LogoUniversiteOffsetY}{%
      \includegraphics[height=\LogoUniversiteHeight,keepaspectratio]{\LogoUniversite}}
  \end{minipage}

  \vspace{\CouvertureEspaceApresLogos}

  % Bandeau de l'entreprise d'accueil.
  \noindent
  \begin{minipage}[c]{0.59\textwidth}
    \centering
    {\fontsize{27}{30}\selectfont\rmfamily SEGURIBAT}%
    \hspace{0.35em}{\fontsize{11}{13}\selectfont\itshape Ingénierie financière}
  \end{minipage}%
  \hfill
  \begin{minipage}[c]{0.37\textwidth}
    \centering
    \fcolorbox{ComptaDark}{white}{%
      \parbox[c][0.82cm][c]{0.91\linewidth}{%
        \centering\fontsize{9.1}{10}\selectfont \MentionSeguribat}}
  \end{minipage}

  \vspace{\CouvertureEspaceApresEntreprise}

  \begin{center}
    \resizebox{0.89\textwidth}{!}{\bfseries \NomEcoleComplet}\\[0.55cm]
    {\fontsize{13.5}{16}\selectfont\bfseries\MakeUppercase{\Formation}}

    \vspace{0.9cm}
    {\color{ComptaBlue}\rule{0.82\textwidth}{0.75pt}}

    \vspace{\CouvertureEspaceAvantTitre}
    {\fontsize{17}{20}\selectfont\bfseries \TypeRapport}\\[0.32cm]
    {\fontsize{43}{47}\selectfont\bfseries
      \textcolor{ComptaBlue}{Compta}\textcolor{ComptaGreen}{Flow}}\\[0.35cm]
    {\fontsize{14}{18}\selectfont \SousTitreRapport}

    \vspace{0.72cm}
    {\color{ComptaBlue}\rule{0.82\textwidth}{0.75pt}}
  \end{center}

  \vspace{\CouvertureEspaceAvantInformations}

  % Bloc d'informations : modifier les commandes placées en tête du fichier.
  \noindent
  \begin{minipage}[t]{0.46\textwidth}
    \raggedright
    \textbf{Réalisé par :}\\[0.18cm]
    \Auteur\\[0.55cm]
    \textbf{Encadrant pédagogique :}\\[0.18cm]
    \EncadrantPedagogique
  \end{minipage}%
  \hfill
  \begin{minipage}[t]{0.43\textwidth}
    \raggedright
    \textbf{Entreprise d'accueil :}\\[0.18cm]
    \EntrepriseAccueil\\[0.55cm]
    \textbf{Encadrant professionnel :}\\[0.18cm]
    \EncadrantEntreprise
  \end{minipage}

  \vfill
  \begin{center}
    \textbf{Période :} \PeriodeProjet\\[0.28cm]
    Année universitaire \AnneeUniversitaire
    \ifx\DateSoutenance\empty\else
      \\[0.28cm]\textbf{Date de soutenance :} \DateSoutenance
    \fi
  \end{center}
  \restoregeometry
\end{titlepage}

\hypersetup{pageanchor=true}
\pagenumbering{roman}

% ============================================================
% 1. Remerciements
% ============================================================
\chapter*{Remerciements}
\addcontentsline{toc}{chapter}{Remerciements}

Au terme de mon stage au sein de \textbf{SEGURIBAT}, je remercie sincèrement toutes les personnes qui ont contribué à cette expérience professionnelle et humaine.

J'adresse une reconnaissance particulière à \textbf{M. Faiz Mohammed Er-Raitab, comptable agréé par l'État, conseiller financier et fiscal et mon encadrant professionnel au sein de SEGURIBAT}, pour sa confiance, sa disponibilité et ses explications. La découverte du domaine comptable et la compréhension progressive de ses besoins métier ont nourri ma réflexion et contribué à faire émerger le projet \textbf{ComptaFlow}. Sa supervision m'a permis de confronter mes choix aux réalités du terrain et d'améliorer progressivement la solution.

Je remercie également \textbf{l'ensemble de l'équipe de SEGURIBAT} pour son accueil, sa bienveillance et son esprit d'entraide. Les échanges avec les collaborateurs m'ont aidée à mieux comprendre le traitement quotidien des dossiers et à développer mon autonomie, mes capacités d'analyse et mes compétences techniques.

J'exprime aussi ma gratitude à \textbf{M. Dehaj} pour son accompagnement et ses orientations, ainsi qu'à \textbf{M. Hnini}, chef de la filière, pour son engagement auprès des étudiants. Je remercie plus largement \textbf{l'École Nationale des Sciences Appliquées de Berrechid} et son corps enseignant pour les connaissances et les méthodes qui m'ont permis d'aborder ce stage avec rigueur.

Enfin, je remercie chaleureusement \textbf{ma famille} pour son soutien constant et ses encouragements. À toutes ces personnes, j'adresse mes sincères remerciements pour leur contribution à cette étape importante de ma formation et de mon parcours d'ingénieure.

% ============================================================
% 2. Résumé
% ============================================================
\chapter*{Résumé}
\addcontentsline{toc}{chapter}{Résumé}

La transformation numérique des cabinets comptables soulève un double enjeu : accélérer le traitement d'un volume important de pièces tout en conservant la fiabilité, la traçabilité et le contrôle humain qu'exige la comptabilité. Le projet \textbf{ComptaFlow} répond à cet enjeu par une application web multi-tenant de gestion documentaire et de pré-comptabilité. La solution permet d'importer des factures et relevés bancaires, d'en extraire les données, d'identifier l'entreprise gérée et le tiers, de préparer des pré-écritures comptables, puis de soumettre les résultats à validation.

L'application repose sur un frontend React et TypeScript, une API FastAPI, une base PostgreSQL et un traitement asynchrone assuré par Celery et Redis. Son pipeline combine Groq Vision, l'extraction de texte native des PDF, PaddleOCR, Groq texte et un repli local Regex/Ollama. Les règles comptables restent déterministes : l'intelligence artificielle détermine la nature économique, tandis que Python recherche le compte exact dans le plan comptable propre à l'entreprise. Le périmètre inclut également les chronos, la détection de doublons, le Grand Livre, la Balance, la TVA, le CPC, le Bilan, les contrôles de pré-clôture, les devises Bank Al-Maghrib et le rapprochement bancaire avancé.

La démarche suivie associe analyse du besoin, conception fonctionnelle, informationnelle et applicative, réalisation incrémentale, puis validation par des tests automatisés portant notamment sur les règles comptables, l'isolation des tenants, les états financiers et la gestion bancaire avancée.

\textbf{Mots-clés :} pré-comptabilité, OCR, intelligence artificielle, FastAPI, React, PostgreSQL, multi-tenant, rapprochement bancaire.

% ============================================================
% 3. Abstract
% ============================================================
\begin{otherlanguage}{english}
\chapter*{Abstract}
\addcontentsline{toc}{chapter}{Abstract}

The digital transformation of accounting firms creates a dual challenge: processing large volumes of documents faster while preserving accounting reliability, auditability, and human oversight. \textbf{ComptaFlow} addresses this challenge through a multi-tenant web platform for document management and pre-accounting. The system imports invoices and bank statements, extracts their data, identifies the managed company and the third party, prepares draft accounting entries, and submits the resulting information for professional review.

The application uses a React and TypeScript frontend, a FastAPI backend, PostgreSQL, and asynchronous processing with Celery and Redis. Its pipeline combines Groq Vision, native PDF text extraction, PaddleOCR, Groq text processing, and a local Regex/Ollama fallback. Accounting decisions remain deterministic: artificial intelligence identifies the economic nature, while Python resolves the exact account from each company's own chart of accounts. The scope also covers document registers, duplicate detection, General Ledger, Trial Balance, VAT, income statement, balance sheet, pre-closing controls, Bank Al-Maghrib exchange rates, and advanced bank reconciliation.

The adopted method combines requirements analysis, functional, information, and application design, incremental implementation, and automated validation of accounting rules, tenant isolation, financial statements, and advanced bank-reconciliation features.

\textbf{Keywords:} pre-accounting, OCR, artificial intelligence, FastAPI, React, PostgreSQL, multi-tenancy, bank reconciliation.
\end{otherlanguage}

% ============================================================
% 4. Listes
% ============================================================
\begingroup
\footnotesize
\tableofcontents
\endgroup
\listoffigures
\listoftables

\chapter*{Liste des abréviations}
\addcontentsline{toc}{chapter}{Liste des abréviations}
\TableSetup
\begin{longtable}{|L{3cm}|L{11cm}|}
\toprule
\rowcolor{ComptaBlue}\TableHead{Abréviation} & \TableHead{Signification} \\
\midrule
API & Application Programming Interface, interface de programmation applicative. \\
BAM & Bank Al-Maghrib. \\
BIC & Bank Identifier Code. \\
CGNC & Code Général de Normalisation Comptable marocain. \\
CNSS & Caisse Nationale de Sécurité Sociale. \\
CPC & Compte de Produits et Charges. \\
CRUD & Create, Read, Update, Delete. \\
HT & Montant hors taxes. \\
IA & Intelligence artificielle. \\
IBAN & International Bank Account Number. \\
ICE & Identifiant Commun de l'Entreprise. \\
JWT & JSON Web Token. \\
OCR & Optical Character Recognition, reconnaissance optique de caractères. \\
OD & Opérations diverses. \\
PDF & Portable Document Format. \\
RIB & Relevé d'Identité Bancaire. \\
SPA & Single Page Application. \\
TTC & Montant toutes taxes comprises. \\
TVA & Taxe sur la valeur ajoutée. \\
UML & Unified Modeling Language. \\
UUID & Universally Unique Identifier. \\
\bottomrule
\end{longtable}

\clearpage
\pagenumbering{arabic}

% ============================================================
% 5. Introduction générale
% ============================================================
\chapter{Introduction générale}

\section{Contexte}

Les cabinets comptables reçoivent des documents hétérogènes : factures d'achat et de vente, relevés bancaires, avis fiscaux et sociaux, fichiers PDF natifs ou images numérisées. Leur traitement manuel nécessite la lecture, le classement, la saisie des montants, l'affectation comptable, la détection des anomalies et la validation. La répétition de ces opérations augmente les délais et le risque d'erreur, en particulier lorsque plusieurs dossiers d'entreprises sont gérés simultanément.

ComptaFlow s'inscrit dans ce contexte comme une plateforme d'assistance. Elle automatise les tâches préparatoires mais maintient le comptable au centre de la décision : les résultats incertains sont signalés, les montants extraits sont préservés pour l'audit et aucune écriture ne peut être validée si elle n'est pas équilibrée.

\section{Problématique}

La problématique générale est la suivante :

\begin{quote}
\textit{Comment automatiser le traitement documentaire et la pré-comptabilité d'un cabinet, depuis le dépôt d'une pièce jusqu'aux états comptables, tout en garantissant l'isolation des dossiers, l'intégrité des montants, la traçabilité et la validation humaine ?}
\end{quote}

Cette problématique se décline en plusieurs difficultés : variabilité de la qualité des documents, identification fiable des rôles entreprise/tiers, résolution des comptes exacts propres à chaque dossier, traitement des devises, rapprochements bancaires complexes et cohérence des états produits.

\section{Objectifs}

Les objectifs du projet sont :

\begin{itemize}[leftmargin=*]
  \item centraliser les pièces comptables dans une interface sécurisée ;
  \item extraire les informations utiles sans altérer silencieusement les valeurs sources ;
  \item classifier les documents et déterminer correctement le sens achat/vente ;
  \item préparer des pré-écritures comptables à partir des informations extraites et du plan comptable de chaque entreprise ;
  \item gérer les relevés, comptes bancaires, paiements partiels et allocations multiples ;
  \item produire les journaux, le Grand Livre, la Balance, la TVA, le CPC et le Bilan ;
  \item fournir des contrôles de cohérence et une piste d'audit ;
  \item assurer la séparation et la confidentialité des données entre les différentes entreprises clientes du cabinet.
\end{itemize}

\section{Méthode suivie}

La réalisation de ComptaFlow s'est appuyée sur une \textbf{démarche incrémentale et itérative inspirée de la méthode Scrum}. Au début du projet, le travail a été organisé sous forme de \textit{sprints} afin de structurer les différentes étapes du développement et de définir progressivement les fonctionnalités à réaliser.

Au fur et à mesure de l'avancement du projet, cette organisation a évolué vers une approche plus flexible. L'approfondissement du domaine comptable, l'étude des règles métier et l'analyse des pratiques observées au sein de SEGURIBAT ont permis d'identifier progressivement de nouveaux besoins, des contraintes supplémentaires ainsi que des améliorations à apporter à l'application.

Le développement s'est ainsi déroulé par cycles successifs. Chaque fonctionnalité était d'abord analysée et conçue, puis implémentée et testée. Des versions intermédiaires de l'application étaient ensuite présentées à \textbf{M. Faiz Mohammed Er-Raitab, encadrant professionnel}, afin de recueillir ses remarques et de vérifier la cohérence du fonctionnement avec les besoins réels du cabinet. Ces retours permettaient de corriger certains éléments, d'améliorer les fonctionnalités existantes ou d'en ajouter de nouvelles.

La méthode de travail suivait donc principalement le cycle suivant : \textbf{analyse du besoin, conception, développement, test, présentation, collecte des retours, puis amélioration}. Cette approche a permis de faire évoluer progressivement ComptaFlow en fonction de la compréhension du domaine comptable et des observations réalisées durant le stage.

Sur le plan technique, chaque évolution était contrôlée au niveau du backend et du frontend. Des tests étaient également réalisés afin de vérifier le bon fonctionnement des fonctionnalités développées et d'identifier les éventuelles erreurs avant leur intégration définitive. Pour les situations comptables ambiguës ou nécessitant une interprétation professionnelle, le système privilégie une \textbf{validation humaine} plutôt qu'une décision automatique non fiable.

Cette démarche a permis de conserver une certaine souplesse tout au long du développement et d'adapter progressivement la solution aux besoins identifiés sur le terrain.

\section{Déroulement du stage}

Le stage s'est déroulé du 1\textsuperscript{er} au 30 juillet 2026. La
chronologie suivante présente les activités dominantes par semaine plutôt que
des dates journalières qui ne pourraient pas toutes être établies avec
précision. Elle s'appuie sur la démarche de travail suivie, les livrables du
projet et l'historique Git disponible. Celui-ci commence le 24 juillet et
confirme principalement les développements, corrections et intégrations de la
dernière semaine ; il ne constitue donc pas une preuve exhaustive des trois
premières semaines.

\begin{table}[H]
\centering
\caption{Déroulement synthétique du stage en juillet 2026}
\footnotesize
\TableSetup
\begin{tabularx}{\textwidth}{|L{2.1cm}|L{6.1cm}|Y|}
\toprule
\rowcolor{ComptaBlue}\TableHead{Période} & \TableHead{Activités principales} & \TableHead{Résultat} \\
\midrule
\shortstack[l]{Semaine 1\\1--7 juillet} & Observation du traitement comptable, découverte de Topaze, échanges métier et formulation du problème. & Processus actuel décrit, besoins initiaux et périmètre fonctionnel dégagés. \\
\shortstack[l]{Semaine 2\\8--14 juillet} & Conception fonctionnelle, informationnelle et applicative ; organisation des rôles, des données et des composants. & Modèles de conception et architecture de référence établis. \\
\shortstack[l]{Semaine 3\\15--21 juillet} & Développement incrémental du socle multi-tenant, de la gestion documentaire, du pipeline d'extraction et de la pré-comptabilité. & Première chaîne fonctionnelle intégrée, soumise aux retours métier. \\
\shortstack[l]{Semaine 4\\22--30 juillet} & Enrichissement des fonctions bancaires et comptables, intégration frontend/backend, tests, corrections et préparation du rapport. & Version consolidée et validée ; les commits des 24, 25, 26 et 27 juillet attestent cette phase d'intégration et d'amélioration. \\
\bottomrule
\end{tabularx}
\end{table}

\section{Organisation du rapport}

Après la présentation de l'entreprise d'accueil, le rapport expose le contexte du projet, puis un état de l'art situe ComptaFlow par rapport à la GED et à la pré-comptabilité assistée. Il présente ensuite les besoins et la conception fonctionnelle, informationnelle et applicative. Les chapitres suivants détaillent la réalisation, le pipeline de traitement, les principales fonctionnalités, les tests et les résultats. La conclusion dresse le bilan, les compétences acquises et les perspectives, avant les références et annexes.

% ============================================================
% 6. Présentation de l'entreprise
% ============================================================
\chapter{Présentation de l'entreprise d'accueil}

Dans le cadre de mon stage d'initiation, j'ai intégré \textbf{SEGURIBAT}, une entreprise spécialisée dans l'accompagnement comptable, fiscal et administratif des entreprises. Elle assure la gestion d'un portefeuille important composé de plusieurs centaines d'entreprises clientes appartenant à différents secteurs d'activité et situées dans plusieurs régions du Maroc.

Malgré le volume important de dossiers traités, SEGURIBAT fonctionne avec une équipe interne relativement réduite, au sein de laquelle les différentes tâches comptables, fiscales et administratives sont réparties entre les collaborateurs. Cette organisation m'a permis d'observer concrètement le traitement des pièces comptables, la saisie des opérations, les déclarations fiscales, le suivi des dossiers clients ainsi que l'utilisation d'outils professionnels tels que \textbf{Topaze}.

Ce chapitre présente l'entreprise d'accueil, ses principales activités, son organisation interne ainsi que les outils et technologies utilisés dans le cadre de ses activités.

\section{Fiche d'identité}
\begin{table}[H]
\centering
\caption{Fiche d'identité de l'entreprise d'accueil}
\TableSetup
\begin{tabularx}{\textwidth}{|L{5cm}|Y|}
\toprule
\rowcolor{ComptaBlue}\TableHead{Élément} & \TableHead{Information} \\
\midrule
\textbf{Raison sociale} & SEGURIBAT \\
\textbf{Secteur d'activité} & Comptabilité, fiscalité, accompagnement administratif et conseil aux entreprises \\
\textbf{Adresse} & 261, boulevard de Témara, Hay Moulay Abdellah, Aïn Chock, Casablanca \\
\textbf{Date de création} & 2015 \\
\textbf{Effectif} & Équipe interne relativement réduite au regard du volume important de dossiers traités \\
\textbf{Portefeuille clients} & Plusieurs centaines d'entreprises appartenant à différents secteurs d'activité \\
\textbf{Zone d'intervention} & Plusieurs villes et régions du Maroc \\
\textbf{Site web} & Non renseigné \\
\textbf{Responsable / Encadrant professionnel} & M. Faiz Mohammed Er-Raitab --- Comptable agréé par l'État, conseiller financier et fiscal \\
\bottomrule
\end{tabularx}
\end{table}

\section{Activités}

SEGURIBAT exerce principalement dans les domaines de la \textbf{comptabilité, de la fiscalité et de l'accompagnement administratif des entreprises}. L'entreprise assure le suivi d'un portefeuille important composé de plusieurs centaines d'entreprises clientes appartenant à différents secteurs d'activité et implantées dans plusieurs villes du Maroc.

Ses activités couvrent notamment la \textbf{tenue de la comptabilité}, depuis la réception et le classement des pièces justificatives jusqu'à la saisie et au contrôle des opérations comptables. Les factures d'achat et de vente, les relevés bancaires ainsi que les différentes pièces nécessaires au suivi des dossiers sont traités puis enregistrés dans les outils comptables utilisés par l'entreprise.

SEGURIBAT intervient également dans la \textbf{gestion fiscale}, notamment à travers la préparation et le traitement des déclarations de TVA ainsi que d'autres obligations fiscales liées à l'activité de ses clients. L'entreprise assure aussi certaines démarches administratives et déclaratives auprès des plateformes et organismes concernés.

Ses prestations comprennent par ailleurs la \textbf{gestion de la paie et des déclarations sociales}, notamment celles liées à la CNSS, ainsi que l'accompagnement des entreprises lors de leur création et dans différentes démarches administratives.

L'entreprise réalise également des travaux de \textbf{contrôle comptable, de préparation des bilans, de clôture des exercices et de conseil fiscal}, permettant aux entreprises clientes de suivre leur situation comptable et de respecter leurs obligations réglementaires.

Compte tenu de l'importance du portefeuille géré, l'activité quotidienne implique le traitement d'un volume conséquent de pièces comptables, de déclarations et d'opérations de saisie, de vérification et de classement, réparties entre les différents collaborateurs de l'entreprise.

\section{Organisation et services}

SEGURIBAT fonctionne avec une \textbf{équipe interne relativement réduite} chargée de gérer un portefeuille important d'entreprises clientes. Son organisation repose principalement sur une \textbf{répartition fonctionnelle des tâches}, plutôt que sur une séparation stricte en plusieurs départements indépendants.

Dans le cadre des activités observées, \textbf{M. Faiz}, mon encadrant professionnel, accompagne les collaborateurs dans le traitement des dossiers et intervient dans les opérations nécessitant davantage d'expertise, de contrôle ou de validation.

Les collaborateurs se répartissent les différentes opérations liées au suivi des dossiers clients. Certains interviennent principalement dans la \textbf{saisie comptable}, notamment l'enregistrement des achats, des ventes et des opérations bancaires. D'autres prennent en charge le \textbf{traitement et la vérification de la TVA}, les déclarations fiscales, la paie et les obligations sociales, ainsi que les travaux de contrôle et de préparation de la clôture comptable.

Bien que les responsabilités soient réparties entre les membres de l'équipe, ces différentes activités restent étroitement liées. Les informations issues du classement et de la saisie des pièces comptables servent notamment aux déclarations fiscales, aux opérations de contrôle ainsi qu'à l'établissement des différents états comptables. Cette organisation nécessite donc une coordination régulière entre les collaborateurs afin d'assurer la continuité et la fiabilité du traitement des dossiers.

Le stage a été réalisé principalement \textbf{au sein de l'activité comptable, avec une dimension informatique transversale, sous la supervision de la direction}.

\section{Organigramme}

L'organigramme suivant présente une représentation fonctionnelle simplifiée de SEGURIBAT, établie à partir de l'organisation et de la répartition des tâches observées durant la période de stage. Les pôles indiqués représentent des familles d'activités et non des départements indépendants.

\begin{figure}[H]
\centering
\resizebox{\textwidth}{!}{%
\begin{tikzpicture}[
    node distance=1.8cm and 1.2cm,
    every node/.style={font=\small},
    main/.style={
        rectangle,
        rounded corners=10pt,
        draw=blue!60!black,
        very thick,
        fill=blue!15,
        minimum width=5cm,
        minimum height=1.2cm,
        align=center,
        text centered
    },
    sub/.style={
        rectangle,
        rounded corners=8pt,
        draw=teal!60!black,
        thick,
        fill=teal!10,
        minimum width=4.2cm,
        minimum height=1.4cm,
        align=center,
        text centered
    },
    orgarrow/.style={
        -{Stealth[length=3mm]},
        thick,
        draw=blue!70!black
    }
]
  \node[main] (direction) {Supervision de l'activité};

  \node[sub, below left=2.2cm and 5.5cm of direction] (compta) {Pôle Comptabilité};
  \node[sub, below left=2.2cm and 1.8cm of direction] (fiscalite) {Pôle Fiscalité\\et Déclarations};
  \node[sub, below right=2.2cm and 1.8cm of direction] (social) {Pôle Social\\et Paie};
  \node[sub, below right=2.2cm and 5.5cm of direction] (admin) {Pôle Administratif\\et Conseil};

  \draw[orgarrow] (direction.south) |- (compta.north);
  \draw[orgarrow] (direction.south) |- (fiscalite.north);
  \draw[orgarrow] (direction.south) |- (social.north);
  \draw[orgarrow] (direction.south) |- (admin.north);
\end{tikzpicture}
}
\caption{Organigramme fonctionnel simplifié de SEGURIBAT}
\label{fig:organigramme-seguribat}
\end{figure}

\section{Technologies utilisées par l'entreprise}

Dans le cadre de ses activités comptables, fiscales et administratives, SEGURIBAT s'appuie principalement sur des outils spécialisés permettant d'assurer la saisie, le suivi et le traitement des dossiers de ses entreprises clientes.

Une grande partie des \textbf{factures et pièces comptables est reçue sous format physique}. Ces documents sont ensuite triés, classés et exploités par les collaborateurs avant la saisie des informations nécessaires dans les outils comptables. Les imprimantes et les scanners sont ainsi régulièrement utilisés pour la gestion, la reproduction et, lorsque cela est nécessaire, la numérisation des documents.

L'outil central utilisé par l'entreprise est \textbf{Topaze}, un logiciel de gestion comptable permettant notamment la saisie et la consultation des écritures, la gestion des journaux d'achat, de vente, de banque, de caisse, d'opérations diverses et de paie, ainsi que le suivi du plan comptable.

Topaze intègre également différentes fonctionnalités liées à la \textbf{TVA, aux immobilisations, à la paie, à la gestion commerciale et aux états de synthèse}, ce qui en fait l'un des principaux outils utilisés quotidiennement par les collaborateurs.

En complément, \textbf{Microsoft Excel} est utilisé ponctuellement pour certains traitements ou contrôles de données. Les échanges professionnels s'effectuent notamment à travers \textbf{Gmail} et \textbf{WhatsApp}.

L'entreprise utilise par ailleurs des \textbf{postes de travail sous Windows}, ainsi que des ressources partagées permettant aux collaborateurs d'accéder aux fichiers nécessaires au traitement des différents dossiers clients.

Les principaux outils observés peuvent être résumés comme suit :

\begin{table}[H]
\centering
\caption{Outils et technologies utilisés par SEGURIBAT}
\label{tab:technologies-seguribat}
\TableSetup
\begin{tabularx}{\textwidth}{|L{3.2cm}|L{4.2cm}|Y|}
\toprule
\rowcolor{ComptaBlue}\TableHead{Catégorie} & \TableHead{Outils et moyens} & \TableHead{Usage principal} \\
\midrule
Gestion documentaire & Documents papier, imprimantes et scanners & Réception, classement, impression et numérisation des pièces \\
Logiciel comptable & Topaze & Saisie, journaux comptables, TVA, paie, immobilisations et états de synthèse \\
Bureautique & Microsoft Excel & Contrôles et traitements ponctuels \\
Communication & Gmail et WhatsApp & Échanges professionnels \\
Infrastructure & Postes Windows et ressources partagées & Accès aux logiciels et aux dossiers clients \\
\bottomrule
\end{tabularx}
\end{table}

% ============================================================
% 7. Contexte du projet
% ============================================================
\chapter{Contexte du projet ComptaFlow}

\section{Problème métier}

Le traitement des pièces comptables au sein de SEGURIBAT repose encore sur plusieurs opérations manuelles, depuis l'identification et le classement des documents jusqu'à l'extraction des informations et leur saisie dans le logiciel comptable Topaze. La figure~\ref{fig:processus-actuel-seguribat} présente le processus observé au sein de l'entreprise.

\begin{figure}[H]
\centering
\resizebox{!}{0.78\textheight}{%
\begin{tikzpicture}[
  node distance=0.3cm,
  currentprocess/.style={
    draw=ComptaGray,
    rounded corners=4pt,
    fill=black!2,
    text width=9.3cm,
    minimum height=0.62cm,
    align=center,
    font=\footnotesize,
    inner sep=3pt
  },
  currentterminal/.style={
    draw=ComptaGreen,
    ellipse,
    fill=ComptaLight,
    minimum width=2.7cm,
    minimum height=0.65cm,
    align=center,
    font=\footnotesize\bfseries
  },
  currentarrow/.style={
    -{Stealth[length=2.2mm]},
    thick,
    draw=ComptaDark
  }
]
  \node[currentterminal] (start) {Début};
  \node[currentprocess, below=of start] (receive) {Réception des documents comptables\\PDF, courriels, scans et autres pièces};
  \node[currentprocess, below=of receive] (print) {Impression des documents};
  \node[currentprocess, below=of print] (company) {Identification de l'entreprise concernée};
  \node[currentprocess, below=of company] (year) {Identification de l'année comptable};
  \node[currentprocess, below=of year] (category) {Identification de la catégorie\\achat, vente, banque, CNSS, TVA, etc.};
  \node[currentprocess, below=of category] (filing) {Classement manuel dans le chrono\\entreprise $\rightarrow$ année $\rightarrow$ catégorie};
  \node[currentprocess, below=of filing] (open) {Ouverture du chrono concerné};
  \node[currentprocess, below=of open] (select) {Sélection de la catégorie à traiter};
  \node[currentprocess, below=of select] (read) {Lecture de la pièce comptable};
  \node[currentprocess, below=of read] (extract) {Extraction manuelle des informations\\date, numéro de facture, ICE / IF / RC et montants HT / TVA / TTC};
  \node[currentprocess, below=of extract] (entry) {Saisie manuelle dans Topaze};
  \node[currentprocess, below=of entry] (stamp) {Apposition du cachet « Document saisi »};
  \node[currentterminal, below=of stamp] (end) {Fin};

  \draw[currentarrow] (start) -- (receive);
  \draw[currentarrow] (receive) -- (print);
  \draw[currentarrow] (print) -- (company);
  \draw[currentarrow] (company) -- (year);
  \draw[currentarrow] (year) -- (category);
  \draw[currentarrow] (category) -- (filing);
  \draw[currentarrow] (filing) -- (open);
  \draw[currentarrow] (open) -- (select);
  \draw[currentarrow] (select) -- (read);
  \draw[currentarrow] (read) -- (extract);
  \draw[currentarrow] (extract) -- (entry);
  \draw[currentarrow] (entry) -- (stamp);
  \draw[currentarrow] (stamp) -- (end);
\end{tikzpicture}%
}
\caption{Processus actuel de traitement d'une pièce comptable au sein de SEGURIBAT}
\label{fig:processus-actuel-seguribat}
\end{figure}

La succession de ces opérations mobilise un temps important, en particulier lorsque le nombre de pièces augmente. Elle expose également le traitement à des erreurs de classement ou de transcription, rend les anomalies moins visibles avant la saisie et complique la traçabilité uniforme des opérations entre les différents dossiers clients. Pour les relevés bancaires, le collaborateur doit en outre examiner les mouvements et rechercher les pièces correspondantes avant leur traitement comptable.

\section{Solutions déjà utilisées et limites}

\begin{table}[H]
\centering
\caption{Solutions utilisées par SEGURIBAT et limites observées}
\label{tab:solutions-existantes-seguribat}
\TableSetup
\begin{tabularx}{\textwidth}{|L{3.3cm}|L{5.1cm}|Y|}
\toprule
\rowcolor{ComptaBlue}\TableHead{Solution utilisée} & \TableHead{Utilisation} & \TableHead{Limites observées} \\
\midrule
\textbf{Chronos physiques} & Classer les pièces par entreprise, année et catégorie. & Classement manuel et recherche lente en cas d'anomalie. \\
\textbf{Traitement manuel des pièces} & Identifier l'entreprise, lire la pièce et relever les informations nécessaires. & Processus répétitif, long et exposé aux erreurs de transcription. \\
\textbf{Topaze} & Saisir les écritures et réaliser une grande partie des traitements comptables. & Les informations doivent encore être préparées et saisies manuellement. \\
\textbf{Microsoft Excel} & Effectuer certains calculs complémentaires avant la saisie dans Topaze. & Les résultats doivent être reportés manuellement dans Topaze, avec un risque d'erreur. \\
\textbf{E-mail et PDF} & Recevoir certaines pièces sous format numérique. & Les documents sont souvent imprimés avant d'être classés dans les chronos. \\
\textbf{Imprimante / scanner} & Imprimer les pièces reçues et numériser certains documents. & Ne supprime ni le classement physique ni la saisie manuelle. \\
\bottomrule
\end{tabularx}
\end{table}

Ces solutions permettent d'assurer le traitement comptable, mais le processus reste \textbf{semi-manuel}, notamment pour le classement, la lecture des pièces, certains calculs et la préparation des données avant leur saisie dans Topaze.

\section{Présentation de ComptaFlow}

ComptaFlow est une plateforme web intelligente destinée aux cabinets comptables. Elle relie la gestion documentaire, l'extraction OCR/IA, les règles de pré-comptabilité et les contrôles professionnels. Une pièce suit un état explicite : en attente, en traitement, traitée, validée ou en erreur. Les écritures suivent de leur côté un cycle brouillon, à vérifier, validé ou rejeté.

Le système sépare trois responsabilités : l'IA extrait et propose une nature économique ; les services Python appliquent les règles ; le plan comptable de l'entreprise fournit le compte exact. Cette séparation réduit le risque d'invention et conserve une intervention humaine lorsque plusieurs comptes sont compatibles ou qu'une donnée est incohérente.

\section{Objectifs spécifiques}

\begin{enumerate}[leftmargin=*]
  \item Centraliser et sécuriser les pièces comptables afin de faciliter leur stockage et leur recherche.
  \item Automatiser l'extraction et le classement des documents, avec des mécanismes de repli en cas d'échec.
  \item Préparer des pré-écritures comptables à partir des informations extraites et du plan comptable de l'entreprise, tout en maintenant une validation humaine lorsque nécessaire.
  \item Assurer la séparation et la confidentialité des données entre les différentes entreprises clientes du cabinet.
  \item Gérer les opérations bancaires et les rapprochements, y compris les paiements partiels et multiples.
  \item Produire les principaux journaux et états comptables et assister la préparation de la clôture.
  \item Préserver la traçabilité des données et faciliter leur exploitation dans le flux comptable existant, notamment avec Topaze.
\end{enumerate}

\section{Périmètre et limites du projet}

ComptaFlow assiste la gestion documentaire et la pré-comptabilité en automatisant
les opérations situées en amont du logiciel comptable. Les tableaux suivants
délimitent les fonctions confirmées dans la version documentée et les éléments
qui ne sont pas revendiqués comme réalisés.

\begin{table}[H]
\centering
\caption{Périmètre fonctionnel couvert par ComptaFlow}
\footnotesize
\TableSetup
\begin{tabularx}{\textwidth}{|L{3.5cm}|Y|}
\toprule
\rowcolor{ComptaBlue}\TableHead{Domaine} & \TableHead{Éléments couverts} \\
\midrule
Accès et isolation & Authentification, rôles, séparation par cabinet et isolation des dossiers d'entreprises. \\
Gestion documentaire & Entreprises, documents, import PDF/PNG/JPEG, chronos numériques et détection des doublons binaires. \\
Extraction et identification & Pipeline OCR/IA, identification de l'entreprise et du tiers, classification achat/vente et contrôles de cohérence. \\
Pré-comptabilité & Préparation et validation humaine des pré-écritures, plan comptable propre à l'entreprise et lignes débit/crédit. \\
Banque & Relevés, comptes bancaires, mouvements, paiements partiels ou multiples et rapprochements par allocations. \\
Devises & Conservation des montants d'origine, conversion séparée en MAD et cours Bank Al-Maghrib. \\
États et contrôles & Grand Livre, Balance, TVA, CPC, Bilan et contrôles de pré-clôture à partir des données validées. \\
Suivi et exploitation & Tableau de bord, tâches, rappels, notifications et exports génériques CSV, XLSX et PDF. \\
\bottomrule
\end{tabularx}
\end{table}

\begin{table}[H]
\centering
\caption{Éléments hors périmètre ou limites de la version documentée}
\footnotesize
\TableSetup
\begin{tabularx}{\textwidth}{|L{4.2cm}|Y|}
\toprule
\rowcolor{ComptaBlue}\TableHead{Élément} & \TableHead{Délimitation retenue} \\
\midrule
Topaze & Le remplacement de Topaze est hors périmètre et aucune intégration automatique directe n'a été réalisée pendant le stage. \\
Autonomie comptable et obligations officielles & La comptabilité sans validation humaine est exclue ; la télétransmission fiscale officielle et un moteur complet de paie ne sont pas intégrés. \\
Canaux d'accès & Aucune application mobile native n'a été réalisée ; l'accès repose sur l'application web responsive. \\
Archivage & Le stockage documentaire actuel n'est pas présenté comme un dispositif d'archivage légal ou probant. \\
Exploitation en production & L'hébergement 24/7, la haute disponibilité et une industrialisation complète ne sont pas revendiqués ; ils constituent des conditions d'évolution. \\
\bottomrule
\end{tabularx}
\end{table}

Les résultats ambigus ou sensibles restent soumis à une validation humaine.
Dans un contexte réel, l'usage de services d'intelligence artificielle externes
nécessiterait en outre une politique explicite de confidentialité et de
protection des documents.

% ============================================================
% État de l'art
% ============================================================
\chapter{État de l'art et positionnement de ComptaFlow}

\section{Gestion électronique des documents}

Dans un cabinet comptable, une gestion électronique des documents organise le
cycle de vie opérationnel des pièces reçues. Elle couvre leur acquisition, leur
classement, l'indexation des informations utiles, la recherche et la
traçabilité des opérations réalisées. Le rattachement à une entreprise, à une
période et à une catégorie permet de retrouver un document dans son contexte
comptable, tandis que le contrôle d'accès limite sa consultation aux
utilisateurs autorisés.

Une GED classique apporte ainsi un classement plus structuré que des dossiers
partagés ou des chronos physiques. Elle ne suffit toutefois pas, à elle seule,
à interpréter le sens comptable de la pièce, à contrôler ses montants ou à
préparer une affectation conforme au plan comptable propre au dossier.

\section{Pré-comptabilité assistée}

La pré-comptabilité regroupe les opérations préparatoires qui précèdent la
validation comptable : collecte des pièces, extraction des données, contrôles,
identification de l'entreprise et du tiers, classification achat/vente,
recherche des comptes possibles, préparation des pré-écritures, contrôle de la
TVA, rapprochement bancaire et préparation des exports. Elle vise à réduire les
tâches répétitives sans transférer au système la responsabilité professionnelle
du comptable.

L'assistance doit donc préserver les données sources, rendre visibles les
incertitudes et laisser à un utilisateur habilité la correction, la validation
ou le rejet. Cette distinction est essentielle : une proposition automatisée
reste une pré-écriture tant qu'elle n'a pas satisfait les règles d'intégrité et
le contrôle humain prévus.

\section{OCR, règles métier et intelligence artificielle}

Une stratégie hybride combine plusieurs techniques complémentaires. La couche
texte d'un PDF natif évite un OCR inutile ; un OCR local traite les documents
numérisés ; des règles déterministes appliquent les contraintes connues ; une
IA peut interpréter des mises en page ou des formulations plus variables. Des
mécanismes de repli limitent la dépendance à une méthode unique, mais aucun
résultat ambigu ne doit être validé automatiquement.

Dans ComptaFlow, Groq Vision constitue le chemin principal. En cas d'échec, le
système tente l'extraction native du PDF ou PaddleOCR. Le texte obtenu est
traité par une branche bancaire spécialisée ou par Groq texte ; si ce service
n'est pas exploitable, les expressions régulières, les heuristiques et un
enrichissement local par Ollama assurent le repli\cite{paddleocr,groq,ollama}.
Les contrôles Python et le plan comptable de l'entreprise déterminent
ensuite les décisions comptables vérifiables. La section consacrée au pipeline
détaille cette séquence sans la répéter ici.

\section{Positionnement de ComptaFlow}

ComptaFlow se situe entre une GED classique et un logiciel comptable complet.
Sa valeur porte sur le traitement amont, la préparation comptable, la
traçabilité, l'association OCR/IA et règles métier, ainsi que le rapprochement
bancaire sous contrôle humain.

\begin{table}[H]
\centering
\caption{Positionnement comparatif des approches de traitement documentaire}
\footnotesize
\TableSetup
\begin{tabularx}{\textwidth}{|L{3.1cm}|L{3.7cm}|L{3.7cm}|Y|}
\toprule
\rowcolor{ComptaBlue}\TableHead{Approche} & \TableHead{Atouts} & \TableHead{Limites} & \TableHead{Adéquation au besoin} \\
\midrule
Dossiers partagés / classement manuel & Simplicité, habitudes existantes et faible coût initial. & Recherche lente, saisies répétitives et traçabilité limitée. & Adapté à un faible volume, moins au portefeuille important observé. \\
GED classique & Centralisation, indexation, recherche et gestion des accès. & Interprétation et préparation comptable souvent limitées. & Répond au classement, mais seulement à une partie du processus amont. \\
SaaS comptable fermé & Fonctions intégrées et exploitation immédiate dans son propre environnement. & Personnalisation, transparence et interopérabilité dépendantes de l'éditeur. & Pertinent si les règles, la confidentialité et les formats correspondent au cabinet. \\
ComptaFlow & Pipeline hybride, règles métier, plan par entreprise, banque et validation humaine. & Solution de stage non industrialisée et sans connecteur direct Topaze. & Couvre la préparation située entre la réception des pièces et le logiciel comptable de référence. \\
\bottomrule
\end{tabularx}
\end{table}

\subsection{Justification synthétique des choix techniques}

Cette comparaison ne constitue pas un benchmark de performances. Elle explicite
les raisons architecturales des choix retenus face à des alternatives possibles.

\begin{table}[H]
\centering
\caption{Choix techniques et alternatives possibles}
\footnotesize
\TableSetup
\begin{tabularx}{\textwidth}{|L{2.2cm}|L{3.4cm}|L{3.4cm}|Y|}
\toprule
\rowcolor{ComptaBlue}\TableHead{Domaine} & \TableHead{Choix retenu} & \TableHead{Alternatives possibles} & \TableHead{Justification} \\
\midrule
Frontend & React et TypeScript & Vue ou Angular & Composants réutilisables, typage et cohérence avec l'architecture du projet. \\
Backend & FastAPI & Django REST ou Node.js & Validation Pydantic, API REST et intégration naturelle avec les traitements Python. \\
Base & PostgreSQL & MySQL ou base NoSQL & Transactions, contraintes relationnelles, types JSONB et fiabilité des données comptables. \\
Asynchrone & Celery et Redis & RQ ou traitement synchrone & Découplage des traitements longs, file de tâches et mécanismes de reprise. \\
OCR & PaddleOCR & Tesseract ou OCR cloud & Solution locale conservée comme mécanisme de repli pour les documents numérisés. \\
\bottomrule
\end{tabularx}
\end{table}

Ces choix s'appuient sur les capacités documentées des composants concernés,
notamment FastAPI, PostgreSQL, Celery et Redis
\cite{fastapi,postgresql,celery,redis}, sans prétendre qu'un benchmark comparatif
des alternatives a été réalisé.

% ============================================================
% 8. Besoins
% ============================================================
\chapter{Analyse des besoins}

\section{Acteurs}

Le système distingue les rôles techniques d'administration des rôles liés au fonctionnement du cabinet. Chaque utilisateur dispose ainsi de droits adaptés à ses responsabilités.

\begin{table}[H]
\centering
\caption{Acteurs de ComptaFlow et responsabilités principales}
\label{tab:acteurs-comptaflow}
\TableSetup
\begin{tabularx}{\textwidth}{|L{4.2cm}|Y|}
\toprule
\rowcolor{ComptaBlue}\TableHead{Acteur} & \TableHead{Responsabilités principales} \\
\midrule
\textbf{Super-administrateur} & Assure l'administration technique globale, la configuration initiale et certaines opérations de maintenance de la plateforme. Durant le développement et les tests, ce rôle a été assuré par la développeuse du projet. \\
\textbf{Administrateur du cabinet} & Gère les utilisateurs du cabinet, attribue les rôles et supervise les paramètres liés à son organisation. \\
\textbf{Expert-comptable} & Contrôle les informations comptables, examine les anomalies et intervient dans les validations sensibles. \\
\textbf{Chef de mission} & Supervise le traitement des dossiers et contrôle l'avancement des opérations comptables. \\
\textbf{Collaborateur} & Importe les documents, consulte les chronos, vérifie les informations extraites et réalise les opérations courantes. \\
\textbf{Assistant} & Participe aux tâches de préparation, d'import et de consultation avec des droits plus limités. \\
\bottomrule
\end{tabularx}
\end{table}

Le \textbf{super-administrateur} possède donc un rôle technique global, distinct des utilisateurs opérationnels du cabinet. Les autres profils correspondent à une organisation progressive des responsabilités métier.

L'accès à l'application est protégé par authentification. Après connexion, les droits associés au rôle de l'utilisateur déterminent les fonctionnalités auxquelles il peut accéder et les opérations qu'il peut effectuer.

\section{Besoins fonctionnels}

\TableSetup
\begin{longtable}{|L{1.4cm}|L{5cm}|L{8cm}|}
\caption{Besoins fonctionnels principaux}\label{tab:bf}\\
\toprule
\rowcolor{ComptaBlue}\TableHead{ID} & \TableHead{Besoin} & \TableHead{Description} \\
\midrule
\endfirsthead
\toprule
\rowcolor{ComptaBlue}\TableHead{ID} & \TableHead{Besoin} & \TableHead{Description} \\
\midrule
\endhead
\TableId{BF01} & Authentification et rôles & Permettre la connexion des utilisateurs et limiter les actions selon le rôle attribué. \\
\TableId{BF02} & Isolation des données & Garantir l'isolation des données entre les cabinets et, au sein de chaque cabinet, entre les différents dossiers d'entreprises selon les droits autorisés. \\
\TableId{BF03} & Import documentaire & Importer des factures, relevés et autres pièces autorisées tout en contrôlant leur format, leur taille et les éventuels doublons. \\
\TableId{BF04} & Extraction des données & Extraire automatiquement les informations utiles des documents et prévoir un mécanisme de repli lorsque l'extraction principale échoue. \\
\TableId{BF05} & Identification et classification & Identifier l'entreprise concernée, reconnaître le tiers et déterminer la nature du document, notamment achat ou vente. \\
\TableId{BF06} & Chronos numériques & Classer automatiquement les documents par entreprise, année, mois et catégorie. \\
\TableId{BF07} & Préparation comptable & Générer une pré-écriture à partir des informations extraites et permettre sa correction, sa validation ou son rejet. \\
\TableId{BF08} & Plan comptable & Importer et gérer le plan comptable propre à chaque entreprise afin d'utiliser les comptes appropriés. \\
\TableId{BF09} & Gestion bancaire & Gérer les relevés et comptes bancaires ainsi que les rapprochements simples, partiels, multiples et les virements internes. \\
\TableId{BF10} & Gestion des devises & Conserver le montant et la devise d'origine, récupérer le taux de change correspondant et effectuer séparément la conversion en MAD. \\
\TableId{BF11} & États comptables & Produire le Grand Livre, la Balance, la TVA, le CPC et le Bilan à partir des données comptables validées. \\
\TableId{BF12} & Pré-clôture & Contrôler la cohérence des comptes, signaler les anomalies et préparer les opérations nécessaires à la clôture sous validation humaine. \\
\TableId{BF13} & Pilotage et suivi & Fournir un tableau de bord, des alertes, des notifications, des tâches et des outils de recherche permettant de suivre l'activité. \\
\TableId{BF14} & Export & Exporter les données et registres et préparer des informations exploitables dans le flux de travail avec Topaze. \\
\bottomrule
\end{longtable}

\section{Besoins non fonctionnels}

Les besoins non fonctionnels définissent les exigences de qualité, de sécurité et de fiabilité auxquelles le système doit répondre.

\TableSetup
\begin{longtable}{|L{1.5cm}|L{4.2cm}|L{8.7cm}|}
\caption{Besoins non fonctionnels}\label{tab:bnf}\\
\toprule
\rowcolor{ComptaBlue}\TableHead{ID} & \TableHead{Qualité} & \TableHead{Exigence} \\
\midrule
\endfirsthead
\toprule
\rowcolor{ComptaBlue}\TableHead{ID} & \TableHead{Qualité} & \TableHead{Exigence} \\
\midrule
\endhead
\TableId{BNF01} & Sécurité & Protéger l'accès par authentification, mots de passe sécurisés et contrôle des droits selon le rôle. \\
\TableId{BNF02} & Confidentialité & Protéger les données entre les cabinets et, dans chaque cabinet, isoler les dossiers des entreprises afin d'empêcher tout accès non autorisé. \\
\TableId{BNF03} & Intégrité des données & Conserver les valeurs extraites d'origine et signaler toute incohérence sans modification silencieuse. \\
\TableId{BNF04} & Fiabilité comptable & Refuser les pré-écritures déséquilibrées et ne jamais générer de compte ou de donnée comptable fictive. \\
\TableId{BNF05} & Traçabilité & Conserver les données d'origine, les statuts, les erreurs et les actions de validation afin de permettre le suivi des traitements. \\
\TableId{BNF06} & Résilience & Prévoir des mécanismes de reprise et des méthodes alternatives lorsque certains traitements d'extraction échouent. \\
\TableId{BNF07} & Performance & Exécuter les traitements lourds sans bloquer l'utilisation de l'interface et maintenir des temps de réponse acceptables. \\
\TableId{BNF08} & Maintenabilité & Organiser l'application en composants et services séparés afin de faciliter les corrections et les évolutions. \\
\TableId{BNF09} & Portabilité & Permettre le déploiement de l'application dans différents environnements avec une configuration adaptée. \\
\TableId{BNF10} & Ergonomie & Proposer une interface claire, responsive et simple à utiliser, avec des statuts, filtres et actions facilement identifiables. \\
\bottomrule
\end{longtable}

\section{Diagramme de cas d'utilisation}

Le diagramme de cas d'utilisation présente les principales interactions entre les différents profils utilisateurs et les fonctionnalités du système. Afin d'améliorer la lisibilité, les rôles ayant des responsabilités proches sont regroupés dans une même représentation.

\begin{figure}[H]
\centering
\resizebox{\textwidth}{!}{%
\begin{tikzpicture}[
  uc/.style={
    draw=ComptaBlue!75!black,
    ellipse,
    fill=white,
    text width=3.75cm,
    minimum height=0.86cm,
    align=center,
    font=\scriptsize,
    inner sep=2.5pt
  },
  zoneTitle/.style={font=\scriptsize\bfseries, text=ComptaDark},
  actorLabel/.style={
    align=center,
    font=\scriptsize\bfseries,
    text=ComptaDark,
    fill=white,
    inner sep=1.5pt
  },
  association/.style={draw=ComptaGray!90!black, line width=0.42pt},
  pics/umlactor/.style={code={
    \draw[draw=ComptaDark, line width=0.65pt] (0,0.55) circle (0.17);
    \draw[draw=ComptaDark, line width=0.65pt] (0,0.38) -- (0,-0.14);
    \draw[draw=ComptaDark, line width=0.65pt] (-0.34,0.18) -- (0,0.27) -- (0.34,0.18);
    \draw[draw=ComptaDark, line width=0.65pt] (0,-0.14) -- (-0.28,-0.52);
    \draw[draw=ComptaDark, line width=0.65pt] (0,-0.14) -- (0.28,-0.52);
  }}
]
  % Frontière du système
  \draw[draw=ComptaGray, line width=0.7pt, rounded corners=5pt]
    (2.45,0.55) rectangle (13.55,16.55);
  \node[font=\small\bfseries, fill=white, inner xsep=5pt]
    at (8,16.55) {Système};

  % Zones fonctionnelles
  \draw[draw=ComptaGray!45, fill=ComptaLight!65, rounded corners=4pt]
    (3.05,10.65) rectangle (7.75,15.35);
  \node[zoneTitle] at (5.40,15.05) {Gestion documentaire};

  \draw[draw=ComptaGray!45, fill=blue!3, rounded corners=4pt]
    (8.25,10.05) rectangle (12.95,15.35);
  \node[zoneTitle] at (10.60,15.05) {Administration};

  \draw[draw=ComptaGray!45, fill=blue!3, rounded corners=4pt]
    (3.05,5.55) rectangle (12.95,9.65);
  \node[zoneTitle] at (8,9.35) {Traitement comptable};

  \draw[draw=ComptaGray!45, fill=ComptaLight!65, rounded corners=4pt]
    (3.05,1.05) rectangle (12.95,5.05);
  \node[zoneTitle] at (8,4.75) {Gestion bancaire};

  % Authentification commune
  \node[uc, text width=3.1cm] (login) at (8,15.95) {S'authentifier};

  % Gestion documentaire
  \node[uc] (document) at (5.40,13.95) {Importer et consulter\\un document};
  \node[uc] (correction) at (5.40,12.65) {Corriger ou contrôler\\les données extraites};
  \node[uc] (chronos) at (5.40,11.35) {Consulter les chronos\\et les registres};

  % Administration
  \node[uc] (users) at (10.60,13.95) {Gérer les utilisateurs\\et attribuer les rôles};
  \node[uc] (settings) at (10.60,12.75) {Gérer le plan comptable\\et les comptes bancaires};
  \node[uc] (pilotage) at (10.60,11.55) {Consulter les outils\\de pilotage};
  \node[uc] (technical) at (10.60,10.45) {Administrer techniquement\\la plateforme};

  % Traitement comptable
  \node[uc, text width=4.4cm] (ledger) at (8,8.55) {Consulter le Grand Livre,\\la Balance et la TVA};
  \node[uc, text width=4.4cm] (statements) at (8,7.30) {Consulter le CPC, le Bilan\\et la pré-clôture};
  \node[uc, text width=4.4cm] (preentry) at (8,6.15) {Contrôler, valider ou rejeter\\une pré-écriture};

  % Gestion bancaire
  \node[uc, text width=4.4cm] (statementsBank) at (8,3.85) {Gérer les relevés bancaires};
  \node[uc, text width=4.4cm] (reconcile) at (8,2.35) {Consulter ou effectuer\\les rapprochements};

  % Acteurs UML à gauche
  \pic at (0.95,12.80) {umlactor};
  \node[actorLabel] at (0.95,11.85) {Assistant /\\Collaborateur};
  \coordinate (assistantLink) at (1.70,12.80);

  \pic at (0.95,7.55) {umlactor};
  \node[actorLabel] at (0.95,6.55) {Expert-comptable /\\Chef de mission};
  \coordinate (expertLink) at (1.70,7.55);

  % Acteurs UML à droite ; le rôle technique est séparé visuellement
  \pic at (15.05,10.45) {umlactor};
  \node[actorLabel, draw=ComptaSlate, rounded corners=2pt]
    at (15.05,9.45) {Super-\\administrateur};
  \coordinate (superLink) at (14.30,10.45);

  \pic at (15.05,13.35) {umlactor};
  \node[actorLabel] at (15.05,12.35) {Administrateur\\du cabinet};
  \coordinate (adminLink) at (14.30,13.35);

  % Associations de l'assistant / collaborateur
  \draw[association] (assistantLink) -- (document.west);
  \draw[association] (assistantLink) -- (correction.west);
  \draw[association] (assistantLink) -- (chronos.west);
  \draw[association] (assistantLink) -- (2.80,12.80) |- (statementsBank.west);
  \draw[association] (assistantLink) -- (2.65,12.80) |- (login.west);

  % Associations de l'expert-comptable / chef de mission
  \draw[association] (expertLink) -- (ledger.west);
  \draw[association] (expertLink) -- (statements.west);
  \draw[association] (expertLink) -- (preentry.west);
  \draw[association] (expertLink) -- (2.90,7.55) |- (document.west);
  \draw[association] (expertLink) -- (2.75,7.55) |- (correction.west);
  \draw[association] (expertLink) -- (2.60,7.55) |- (chronos.west);
  \draw[association] (expertLink) -- (2.45,7.55) |- (reconcile.west);
  \draw[association] (expertLink) -- (2.30,7.55) |- (login.west);

  % Associations de l'administrateur du cabinet
  \draw[association] (adminLink) -- (users.east);
  \draw[association] (adminLink) -- (settings.east);
  \draw[association] (adminLink) -- (pilotage.east);
  \draw[association] (adminLink) -- (13.15,13.35) |- (login.east);

  % Associations du super-administrateur technique
  \draw[association] (superLink) -- (technical.east);
  \draw[association] (superLink) -- (13.35,10.45) |- (users.east);
  \draw[association] (superLink) -- (13.50,10.45) |- (login.east);
\end{tikzpicture}%
}
\caption{Diagramme général des cas d'utilisation}
\label{fig:usecase}
\end{figure}

Les assistants et collaborateurs interviennent principalement dans le traitement courant des documents. Les chefs de mission et experts-comptables disposent de fonctions supplémentaires de contrôle, de consultation des états comptables et de validation. L'administrateur du cabinet assure principalement la gestion des utilisateurs et des paramètres du cabinet, tandis que le super-administrateur intervient au niveau de l'administration technique globale.

% ============================================================
% 9. Conception
% ============================================================
\chapter{Conception du système}

\section{Conception fonctionnelle}

Le traitement principal d'une pièce comptable suit un enchaînement allant de son import jusqu'au contrôle et à la validation des données obtenues.

\subsection{Scénario nominal d'une facture}

\begin{enumerate}[leftmargin=*]
  \item L'utilisateur importe une facture au format PDF ou image.
  \item Le système vérifie le document et l'enregistre pour traitement.
  \item Les informations utiles de la pièce sont extraites automatiquement.
  \item Le système identifie l'entreprise concernée, le tiers et détermine s'il s'agit notamment d'un achat ou d'une vente.
  \item Le document est classé dans le chrono correspondant selon l'entreprise, la période et la catégorie.
  \item À partir des informations extraites et du plan comptable de l'entreprise, une pré-écriture comptable est préparée.
  \item L'utilisateur contrôle la pièce, les données extraites et les comptes proposés. Il peut corriger les informations si nécessaire.
  \item Un utilisateur autorisé valide ou rejette la pré-écriture.
  \item Les données validées peuvent ensuite alimenter le Grand Livre, la Balance et les autres états comptables.
\end{enumerate}

\subsection{Diagramme d'activité}

\begin{figure}[H]
\centering
\resizebox{!}{0.72\textheight}{%
\begin{tikzpicture}[node distance=0.72cm]
  \node[box, fill=ComptaLight] (start) {Sélection du document};
  \node[box, below=of start] (valid) {Contrôles fichier et doublon};
  \node[box, below=of valid] (queue) {Enregistrement du document\\pour traitement};
  \node[box, below=of queue] (extract) {Extraction et classification};
  \node[box, below=of extract] (company) {Identification entreprise / tiers};
  \node[diamond, draw, below=0.85cm of company, aspect=2.2, align=center, font=\small] (type) {Document bancaire ?};
  \node[box, below left=0.9cm and 1.4cm of type] (movement) {Traiter les mouvements et\\rapprochements bancaires};
  \node[box, below right=0.9cm and 1.4cm of type] (entry) {Préparer la pré-écriture\\comptable};
  \node[diamond, draw, below=2.2cm of type, aspect=2.2, align=center, font=\small] (check) {Données cohérentes\\et suffisantes ?};
  \node[box, below left=0.9cm and 1.4cm of check] (manual) {Vérification / correction\\humaine};
  \node[box, below right=0.9cm and 1.4cm of check] (proposal) {Soumettre à validation};
  \node[diamond, draw, below=2.2cm of check, aspect=2.2, align=center, font=\small] (approval) {Pré-écriture\\validée ?};
  \node[box, fill=ComptaLight, below=0.9cm of approval, minimum width=6.2cm] (states) {Alimentation du Grand Livre, de la Balance\\et des états comptables};

  \draw[flow] (start) -- (valid);
  \draw[flow] (valid) -- (queue);
  \draw[flow] (queue) -- (extract);
  \draw[flow] (extract) -- (company);
  \draw[flow] (company) -- (type);
  \draw[flow] (type) -| node[above left,font=\scriptsize]{oui} (movement);
  \draw[flow] (type) -| node[above right,font=\scriptsize]{non} (entry);
  \draw[flow] (movement) |- (check);
  \draw[flow] (entry) |- (check);
  \draw[flow] (check) -| node[above left,font=\scriptsize]{non} (manual);
  \draw[flow] (check) -| node[above right,font=\scriptsize]{oui} (proposal);
  \draw[flow] (manual.east) -- node[above,font=\scriptsize]{après correction} (proposal.west);
  \draw[flow] (proposal.south) |- (approval.east);
  \draw[flow] (approval.west) -| node[near start,below,font=\scriptsize]{non} ([xshift=-0.7cm]manual.west) -- (manual.west);
  \draw[flow] (approval) -- node[right,font=\scriptsize]{oui} (states);
\end{tikzpicture}
}
\caption{Activité principale de traitement d'un document}
\label{fig:activity}
\end{figure}

\clearpage
\subsection{Diagramme de séquence du traitement d'une pièce}

Afin de compléter la représentation fonctionnelle précédente, le diagramme suivant illustre les échanges entre les principaux composants. Il met en évidence le traitement asynchrone, l'enregistrement progressif des résultats et le contrôle humain avant validation.

\begin{figure}[H]
\centering
\resizebox{0.94\textwidth}{!}{%
\begin{tikzpicture}[
  x=1cm,
  y=1cm,
  seqParticipant/.style={
    draw=ComptaBlue,
    rounded corners=3pt,
    fill=blue!5,
    minimum width=2.05cm,
    minimum height=0.8cm,
    text width=1.9cm,
    align=center,
    font=\scriptsize\bfseries
  },
  seqLife/.style={draw=ComptaGray, dashed, thin},
  seqMessage/.style={-{Latex[length=2mm]}, draw=ComptaDark, thick},
  seqReturn/.style={-{Latex[length=2mm]}, draw=ComptaGray, dashed},
  seqLabel/.style={font=\scriptsize, fill=white, inner sep=1.5pt, align=center},
  seqNote/.style={
    draw=ComptaSlate,
    rounded corners=3pt,
    fill=ComptaGray!6,
    text width=4.2cm,
    align=left,
    font=\scriptsize
  }
]
  % Participants
  \node[seqParticipant] (user) at (0,0) {Utilisateur};
  \node[seqParticipant] (web) at (2.35,0) {Interface Web};
  \node[seqParticipant] (api) at (4.7,0) {API FastAPI};
  \node[seqParticipant] (db) at (7.05,0) {Base PostgreSQL};
  \node[seqParticipant] (redis) at (9.4,0) {Redis / file de tâches};
  \node[seqParticipant] (worker) at (11.75,0) {Worker Celery};
  \node[seqParticipant] (services) at (14.1,0) {Services OCR/IA et métier};

  % Lignes de vie
  \foreach \participant in {user,web,api,db,redis,worker,services}{
    \draw[seqLife] (\participant.south) -- ++(0,-15.85);
  }

  % Import et contrôles initiaux
  \draw[seqMessage] (0,-0.85) -- node[seqLabel,above]{Sélectionner une pièce} (2.35,-0.85);
  \draw[seqMessage] (2.35,-1.45) -- node[seqLabel,above]{Importer le fichier} (4.7,-1.45);
  \draw[seqMessage] (4.7,-2.0) -- ++(0.72,0) -- ++(0,-0.5) -- (4.7,-2.5);
  \node[seqLabel, font=\tiny, anchor=west] at (5.43,-2.25) {Format / taille / doublon / stockage local};
  \draw[seqMessage] (4.7,-3.05) -- node[seqLabel,above]{Créer le document\\en attente} (7.05,-3.05);
  \draw[seqMessage] (4.7,-3.65) -- node[seqLabel,above]{Publier \texttt{process\_document(id)}} (9.4,-3.65);
  \draw[seqReturn] (4.7,-4.25) -- node[seqLabel,above]{Document créé} (2.35,-4.25);

  % Traitement asynchrone
  \node[font=\scriptsize\bfseries, text=ComptaGreen, fill=white] at (10.55,-4.7) {Traitement asynchrone};
  \draw[seqMessage] (9.4,-5.15) -- node[seqLabel,above]{Distribuer la tâche} (11.75,-5.15);
  \draw[seqMessage] (11.75,-5.75) -- node[seqLabel,above]{Statut : en traitement} (7.05,-5.75);
  \draw[seqMessage] (11.75,-6.35) -- node[seqLabel,above]{Extraire et classifier} (14.1,-6.35);
  \draw[seqReturn] (14.1,-6.95) -- node[seqLabel,above]{Texte et données extraites} (11.75,-6.95);
  \node[seqLabel, font=\scriptsize\itshape, text=ComptaSlate, anchor=east] at (14.1,-7.45) {Repli si l'extraction est insuffisante};
  \draw[seqMessage] (11.75,-8.55) -- node[seqLabel,above]{Identifier, classer et préparer\\la pré-écriture} (14.1,-8.55);
  \draw[seqReturn] (14.1,-9.15) -- node[seqLabel,above]{Résultat et anomalies éventuelles} (11.75,-9.15);
  \draw[seqMessage] (11.75,-9.75) -- node[seqLabel,above]{Enregistrer résultats, statut\\et pré-écriture} (7.05,-9.75);
  \node[seqLabel, font=\scriptsize\itshape, text=ComptaSlate, anchor=east] at (14.1,-10.35) {Ambiguïté : données conservées, contrôle humain requis};

  % Consultation puis décision humaine
  \draw[seqMessage] (2.35,-11.25) -- node[seqLabel,above]{Consulter l'avancement} (4.7,-11.25);
  \draw[seqMessage] (4.7,-11.85) -- node[seqLabel,above]{Lire statut et résultat} (7.05,-11.85);
  \draw[seqReturn] (7.05,-12.45) -- node[seqLabel,above]{Données persistées} (4.7,-12.45);
  \draw[seqReturn] (4.7,-13.05) -- node[seqLabel,above]{Afficher pièce et pré-écriture} (2.35,-13.05);
  \draw[seqMessage] (0,-13.65) -- node[seqLabel,above]{Contrôler et, si nécessaire, corriger} (2.35,-13.65);
  \draw[seqMessage] (2.35,-14.25) -- node[seqLabel,above]{Correction ou validation / rejet} (4.7,-14.25);
  \draw[seqMessage] (4.7,-14.85) -- node[seqLabel,above]{Mettre à jour la pré-écriture} (7.05,-14.85);
  \draw[seqReturn] (4.7,-15.45) -- node[seqLabel,above]{Confirmation} (2.35,-15.45);
  \draw[seqReturn] (2.35,-16.05) -- node[seqLabel,above]{Résultat de la décision} (0,-16.05);
\end{tikzpicture}%
}
\caption{Diagramme de séquence du traitement d'une pièce comptable}
\label{fig:sequence-document}
\end{figure}

L'interface transmet le document sans exécuter elle-même les traitements lourds. L'API contrôle l'import, organise la persistance et déclenche le traitement asynchrone par l'intermédiaire de la file de tâches. Le worker coordonne ensuite l'extraction et l'analyse métier, puis enregistre progressivement les résultats et les anomalies éventuelles. La pré-écriture ainsi préparée reste soumise au contrôle, à la correction et à la décision d'un utilisateur autorisé avant son exploitation dans les fonctions comptables sensibles.

\subsection{Règles fonctionnelles critiques}

Certaines décisions nécessitent des règles métier précises afin d'éviter une classification ou un traitement comptable incorrect.

\begin{table}[H]
\centering
\caption{Principales règles fonctionnelles appliquées}
\label{tab:regles-fonctionnelles-critiques}
\TableSetup
\begin{tabularx}{\textwidth}{|L{4.1cm}|Y|}
\toprule
\rowcolor{ComptaBlue}\TableHead{Règle} & \TableHead{Principe appliqué} \\
\midrule
Identification achat / vente & Pour un achat, l'entreprise gérée est destinataire et le tiers est fournisseur. Pour une vente, l'entreprise est émettrice et le tiers est client. La classification ne repose donc pas uniquement sur le nom du tiers. \\
Contrôle des montants & Les montants HT, TVA et TTC extraits de la pièce sont conservés comme données de référence. Les calculs réalisés servent principalement à vérifier leur cohérence et à signaler les anomalies. \\
Sélection des comptes & Les comptes sont recherchés dans le plan comptable propre à l'entreprise. Lorsque plusieurs comptes sont possibles ou que le choix est ambigu, le système demande une vérification humaine au lieu de sélectionner arbitrairement un compte. \\
\bottomrule
\end{tabularx}
\end{table}

\section{Conception informationnelle}

\subsection{Principes de modélisation}

Les entités utilisent des UUID et des horodatages. La majorité des données métier propres au cabinet sont rattachées à un \texttt{cabinet\_id}, tandis que les données propres à un dossier comptable sont également associées à un \texttt{entreprise\_id}. Cette redondance contrôlée de portée facilite les filtres de sécurité et empêche qu'un identifiant métier seul donne accès à un autre tenant.

Dans l'implémentation actuelle, le chrono constitue le conteneur documentaire principal d'une entreprise. Il n'est pas créé un chrono distinct pour chaque exercice ou catégorie : le classement interne est porté par les attributs \texttt{annee}, \texttt{mois} et \texttt{categorie} du document. La contrainte \texttt{UNIQUE(entreprise\_id)} de la table \texttt{chronos} traduit donc bien un maximum d'un chrono principal par entreprise.

\subsection{Modèle conceptuel de données Merise}

Le MCD exprime les objets métier et leurs associations sans dépendre de PostgreSQL. Les cardinalités sont écrites dans le sens entité--association. Par exemple, un cabinet gère de zéro à plusieurs entreprises, tandis qu'une entreprise appartient obligatoirement à un seul cabinet. Un document peut rester sans entreprise et sans chrono lorsqu'il concerne directement le cabinet ; cette possibilité est effectivement prévue par le modèle et par le traitement documentaire.

\begin{figure}[H]
\centering
\resizebox{\textwidth}{!}{%
\begin{tikzpicture}[x=1cm,y=1cm]
  \node[meriseEntity] (cabinet) at (0,0) {\textbf{CABINET}\nodepart{second}\underline{id\_cabinet}\\nom\\ICE\\coordonnées\\actif};
  \node[meriseAssoc] (appartenir) at (3,2.25) {APPARTENIR};
  \node[meriseEntity] (user) at (6,2.25) {\textbf{UTILISATEUR}\nodepart{second}\underline{id\_utilisateur}\\email\\nom, prénom\\rôle\\actif};
  \node[meriseAssoc] (gerer) at (3,-2.2) {GÉRER};
  \node[meriseEntity] (company) at (6,-2.2) {\textbf{ENTREPRISE}\nodepart{second}\underline{id\_entreprise}\\nom\\ICE, IF, RC\\forme juridique\\actif};

  \node[meriseAssoc] (deposer) at (9,1.45) {DÉPOSER};
  \node[meriseEntity] (document) at (12,1.45) {\textbf{DOCUMENT}\nodepart{second}\underline{id\_document}\\fichier, hash\\année, mois, catégorie\\statut, texte OCR\\données extraites};
  \node[meriseAssoc] (posseder) at (9,-3.65) {POSSÉDER};
  \node[meriseEntity] (chrono) at (12,-3.65) {\textbf{CHRONO}\nodepart{second}\underline{id\_chrono}\\nom};
  \node[meriseAssoc] (classer) at (12,-1.15) {CLASSER};
  \node[meriseAssoc] (rattacher) at (9,-0.35) {RATTACHER};

  \node[meriseAssoc] (produire) at (15,2.45) {PRODUIRE};
  \node[meriseEntity] (entry) at (18,2.45) {\textbf{PRÉ-ÉCRITURE COMPTABLE}\nodepart{second}\underline{id\_pré-écriture}\\type, date, tiers\\HT, TVA, TTC\\devise\\statut validation};
  \node[meriseAssoc] (extraire) at (15,-1.05) {EXTRAIRE};
  \node[meriseEntity] (movement) at (18,-1.05) {\textbf{MOUVEMENT}\nodepart{second}\underline{id\_mouvement}\\date, libellé\\débit/crédit\\montant, devise\\statut rapprochement};
  \node[meriseAssoc] (rapprocher) at (21.2,0.7) {RAPPROCHER\\montant affecté\\statut, score};

  \node[meriseAssoc] (planifier) at (3,-5.25) {PLANIFIER};
  \node[meriseEntity] (account) at (0,-5.25) {\textbf{COMPTE COMPTABLE}\nodepart{second}\underline{id\_compte}\\numéro, libellé\\famille CGNC\\usage, nature};
  \node[meriseAssoc] (detenir) at (7.8,-5.25) {DÉTENIR};
  \node[meriseEntity] (bankaccount) at (11,-6.15) {\textbf{COMPTE BANCAIRE}\nodepart{second}\underline{id\_compte bancaire}\\banque\\RIB, IBAN, BIC\\devise\\compte comptable};
  \node[meriseAssoc] (affecter) at (15,-4.75) {AFFECTER};

  \draw (cabinet) -- node[pos=.22,above left,font=\tiny]{0,N} (appartenir);
  \draw (appartenir) -- node[pos=.78,above,font=\tiny]{1,1} (user);
  \draw (cabinet) -- node[pos=.2,below left,font=\tiny]{0,N} (gerer);
  \draw (gerer) -- node[pos=.78,below,font=\tiny]{1,1} (company);
  \draw (user) -- node[pos=.2,above,font=\tiny]{0,N} (deposer);
  \draw (deposer) -- node[pos=.78,above,font=\tiny]{1,1} (document);
  \draw (company) -- node[pos=.2,below,font=\tiny]{0,1} (posseder);
  \draw (posseder) -- node[pos=.78,below,font=\tiny]{1,1} (chrono);
  \draw (chrono) -- node[pos=.18,right,font=\tiny]{0,N} (classer);
  \draw (classer) -- node[pos=.78,right,font=\tiny]{0,1} (document);
  \draw (company) -- node[pos=.2,above,font=\tiny]{0,N} (rattacher);
  \draw (rattacher) -- node[pos=.78,above,font=\tiny]{0,1} (document);
  \draw (document) -- node[pos=.2,above,font=\tiny]{0,1} (produire);
  \draw (produire) -- node[pos=.78,above,font=\tiny]{1,1} (entry);
  \draw (document) -- node[pos=.2,below,font=\tiny]{0,N} (extraire);
  \draw (extraire) -- node[pos=.78,below,font=\tiny]{1,1} (movement);
  \draw (entry) -- node[pos=.22,above,font=\tiny]{0,N} (rapprocher);
  \draw (movement) -- node[pos=.22,below,font=\tiny]{0,N} (rapprocher);
  \draw (account) -- node[pos=.2,above,font=\tiny]{1,1} (planifier);
  \draw (planifier) -- node[pos=.78,above,font=\tiny]{0,N} (company);
  \draw (company) -- node[pos=.22,below,font=\tiny]{0,N} (detenir);
  \draw (detenir) -- node[pos=.78,above,font=\tiny]{1,1} (bankaccount);
  \draw (bankaccount) -- node[pos=.2,above,font=\tiny]{0,N} (affecter);
  \draw (affecter) -- node[pos=.78,below,font=\tiny]{0,1} (movement);
\end{tikzpicture}%
}
\caption{MCD Merise du cœur documentaire, comptable et bancaire}
\label{fig:mcd-merise}
\end{figure}

Le rapprochement est une association porteuse de propriétés. Une pré-écriture comptable peut être réglée par plusieurs mouvements et un mouvement peut être réparti entre plusieurs pré-écritures. Le montant affecté, le score et le statut appartiennent donc à l'association \textsc{Rapprocher}, et non à l'une des deux entités. Cette modélisation correspond à la table d'allocation utilisée pour les paiements partiels et multiples.

\begin{table}[H]
\centering
\caption{Règles de cardinalité du MCD}
\TableSetup
\begin{tabularx}{\textwidth}{|L{3.5cm}|L{3cm}|Y|}
\toprule
\rowcolor{ComptaBlue}\TableHead{Association} & \TableHead{Cardinalités} & \TableHead{Règle métier} \\
\midrule
APPARTENIR & Cabinet 0,N ; utilisateur 1,1 & Tout utilisateur dépend d'un seul cabinet. \\
GÉRER & Cabinet 0,N ; entreprise 1,1 & Une entreprise gérée ne peut appartenir qu'à un cabinet. \\
DÉPOSER & Utilisateur 0,N ; document 1,1 & Chaque document conserve l'auteur de son dépôt. \\
POSSÉDER & Entreprise 0,1 ; chrono 1,1 & Une entreprise possède au plus un chrono principal ; année, mois et catégorie sont portés par les documents. \\
CLASSER & Chrono 0,N ; document 0,1 & Un document du cabinet peut rester hors chrono entreprise. \\
RATTACHER & Entreprise 0,N ; document 0,1 & Un document peut rester au niveau du cabinet sans dossier artificiel. \\
PRODUIRE & Document 0,1 ; pré-écriture 1,1 & Une facture produit au plus une pré-écriture, composée ensuite de plusieurs lignes débit/crédit. \\
EXTRAIRE & Document 0,N ; mouvement 1,1 & Chaque mouvement extrait garde son document source. \\
RAPPROCHER & Pré-écriture 0,N ; mouvement 0,N & L'association couvre paiements partiels et allocations multiples. \\
PLANIFIER & Entreprise 0,N ; compte 1,1 & Tout compte exact appartient au plan d'une seule entreprise. \\
DÉTENIR & Entreprise 0,N ; compte bancaire 1,1 & Une entreprise peut gérer plusieurs banques. \\
AFFECTER & Compte bancaire 0,N ; mouvement 0,1 & L'affectation reste facultative tant que le RIB/IBAN n'est pas résolu. \\
\bottomrule
\end{tabularx}
\end{table}

\subsection{Modèle logique de données Merise}

Le passage au MLD transforme les associations de type 1,N en clés étrangères placées du côté N. L'association N,N \textsc{Rapprocher} devient une relation autonome. La notation \textbf{PK} indique une clé primaire, \textit{FK} une clé étrangère et \textsc{UQ} une contrainte d'unicité.

\small
\TableSetup
\begin{longtable}{|L{4.4cm}|L{10cm}|}
\caption{MLD relationnel de ComptaFlow}\label{tab:mld-merise}\\
\toprule
\rowcolor{ComptaBlue}\TableHead{Relation} & \TableHead{Attributs logiques principaux} \\
\midrule
\endfirsthead
\toprule
\rowcolor{ComptaBlue}\TableHead{Relation} & \TableHead{Attributs logiques principaux} \\
\midrule
\endhead
CABINET & \textbf{PK} id, nom, ICE \textsc{UQ}, adresse, téléphone, email, actif \\
UTILISATEUR & \textbf{PK} id, \textit{FK} cabinet\_id, email \textsc{UQ}, mot\_de\_passe\_haché, nom, prénom, rôle, actif \\
ENTREPRISE & \textbf{PK} id, \textit{FK} cabinet\_id, nom, ICE, IF, RC, forme\_juridique, adresse, actif ; \textsc{UQ}(cabinet\_id, ICE) \\
CHRONO & \textbf{PK} id, \textit{FK} cabinet\_id, \textit{FK} entreprise\_id \textsc{UQ}, nom \\
DOCUMENT & \textbf{PK} id, \textit{FK} cabinet\_id, \textit{FK?} entreprise\_id, \textit{FK?} chrono\_id, \textit{FK} uploaded\_by, année, mois, catégorie, fichier, hash, statut, texte\_ocr, données\_extraites \\
\shortstack[l]{PRÉ-ÉCRITURE\\COMPTABLE\\(\texttt{ecritures\_}\\\texttt{comptables})} & \textbf{PK} id, \textit{FK} cabinet\_id, \textit{FK} entreprise\_id, \textit{FK} document\_id, \textit{FK?} validated\_by, type, comptes, HT, TVA, TTC, devise, statut, dates et référence de saisie Topaze \\
LIGNE COMPTABLE & \textbf{PK} id, \textit{FK} cabinet\_id, \textit{FK} entreprise\_id, \textit{FK?} écriture\_id, \textit{FK?} mouvement\_id, \textit{FK?} régularisation\_id, \textit{FK?} période\_TVA\_id, date, journal, compte, débit, crédit, ordre, validée \\
\shortstack[l]{COMPTE COMPTABLE\\ENTREPRISE} & \textbf{PK} id, \textit{FK} cabinet\_id, \textit{FK} entreprise\_id, numéro, libellé, famille, usage, nature, tiers ; \textsc{UQ}(cabinet\_id, entreprise\_id, numéro) \\
\shortstack[l]{COMPTE BANCAIRE\\ENTREPRISE} & \textbf{PK} id, \textit{FK} cabinet\_id, \textit{FK} entreprise\_id, libellé, banque, RIB, IBAN, BIC, devise, numéro\_compte\_comptable \\
MOUVEMENT BANCAIRE & \textbf{PK} id, \textit{FK} cabinet\_id, \textit{FK} entreprise\_id, \textit{FK} document\_id, \textit{FK?} compte\_bancaire\_id, \textit{FK?} mouvement\_lié\_id, date, libellé, type\_mouvement, montant, devise, rapprochement \\
\shortstack[l]{RAPPROCHEMENTS\\BANCAIRES\\ALLOCATIONS} & \textbf{PK} id, \textit{FK} cabinet\_id, \textit{FK} entreprise\_id, \textit{FK} mouvement\_bancaire\_id, \textit{FK} écriture\_id, montant\_affecté, écart\_change, statut, score ; \textsc{UQ}(mouvement\_bancaire\_id, écriture\_id) \\
\shortstack[l]{TVA\\CONFIGURATIONS\\ENTREPRISE} & \textbf{PK} id, \textit{FK} cabinet\_id, \textit{FK} entreprise\_id, périodicité, prorata, retenue, comptes exacts de TVA collectée, récupérable, à payer et crédit ; \textsc{UQ}(cabinet\_id, entreprise\_id) \\
TVA\_PÉRIODE & \textbf{PK} id, \textit{FK} cabinet\_id, \textit{FK} entreprise\_id, année, mois, statut comptable, statut déclaratif, collectée, récupérable, crédits, TVA\_nette, échéance, déclaration et référence Topaze ; \textsc{UQ}(cabinet\_id, entreprise\_id, année, mois) \\
TVA RÉGULARISATION & \textbf{PK} id, \textit{FK} cabinet\_id, \textit{FK} entreprise\_id, \textit{FK} période\_id, nature, montant\_signé, motif, date, statut \\
\shortstack[l]{TVA CRÉDIT\\UTILISATION} & \textbf{PK} id, \textit{FK} entreprise\_id, \textit{FK} période\_source\_id \textsc{UQ}, \textit{FK} période\_destination\_id \textsc{UQ}, montant\_utilisé \\
\shortstack[l]{RÉGULARISATION\\CLÔTURE} & \textbf{PK} id, \textit{FK} cabinet\_id, \textit{FK} entreprise\_id, exercice, date, type, libellé, montant, comptes, statut, anomalies \\
TAUX CHANGE BAM & \textbf{PK} id, date\_cours, devise, unité, cours\_achat, cours\_vente, source ; \textsc{UQ}(date\_cours, devise, source) \\
ANOMALIE COMPTABLE & \textbf{PK} id, \textit{FK} cabinet\_id, \textit{FK} entreprise\_id, \textit{FK?} écriture\_id, code, niveau, message, statut, résolution \\
DOCUMENT ATTENDU & \textbf{PK} id, \textit{FK} cabinet\_id, \textit{FK} entreprise\_id, type, fréquence, période, échéance, quantité attendue et reçue, statut \\
PÉRIODE DE TRAVAIL & \textbf{PK} id, \textit{FK} cabinet\_id, \textit{FK} entreprise\_id, exercice, début, fin, statut, verrouillage, réouverture et justification \\
TÂCHE & \textbf{PK} id, \textit{FK} cabinet\_id, \textit{FK?} entreprise\_id, \textit{FK} cree\_par, \textit{FK?} assignee\_a, titre, date et heure d'échéance, priorité, récurrence, statut \\
MESSAGE CABINET & \textbf{PK} id, \textit{FK} cabinet\_id, \textit{FK} expéditeur\_id, \textit{FK} destinataire\_id, contenu, lecture, date d'envoi \\
AUDIT LOG & \textbf{PK} id, \textit{FK} cabinet\_id, \textit{FK?} entreprise\_id, \textit{FK?} user\_id, action, module, ressource, valeurs avant/après, métadonnées, adresse IP et horodatage \\
\bottomrule
\end{longtable}
\normalsize

Le MCD retient la règle fonctionnelle d'une seule pré-écriture par facture. Le service de préparation comptable applique cette règle en réutilisant la première pré-écriture trouvée et en supprimant les anciens doublons. En revanche, dans le schéma physique actuel, \texttt{ecritures\_comptables.document\_id} est une clé étrangère indexée mais non unique. Le MLD ne lui attribue donc pas de contrainte \textsc{UQ} inexistante ; une future migration pourrait renforcer cette règle au niveau de la base après contrôle des données existantes.

\subsection{Modèle physique de données PostgreSQL}

Le MPD applique le MLD à PostgreSQL. Les identifiants sont des \texttt{UUID}, les montants utilisent \texttt{NUMERIC} afin d'éviter les erreurs d'arrondi binaires, et les données d'extraction auditables sont stockées en \texttt{JSONB}. Les objets dépendants utilisent des actions référentielles explicites, principalement \texttt{ON DELETE CASCADE} pour les portées cabinet et entreprise et \texttt{SET NULL} pour certaines références optionnelles.

\small
\TableSetup
\begin{longtable}{|L{4cm}|L{4.2cm}|L{6.1cm}|}
\caption{MPD PostgreSQL : types et contraintes structurantes}\label{tab:mpd-merise}\\
\toprule
\rowcolor{ComptaBlue}\TableHead{Table / colonne} & \TableHead{Type physique} & \TableHead{Contraintes ou usage} \\
\midrule
\endfirsthead
\toprule
\rowcolor{ComptaBlue}\TableHead{Table / colonne} & \TableHead{Type physique} & \TableHead{Contraintes ou usage} \\
\midrule
\endhead
Toutes les tables métier : \texttt{id} & \texttt{UUID} & Clé primaire, valeur générée par UUID4. \\
Tables multi-tenant : \texttt{cabinet\_id} & \texttt{UUID} & FK vers \texttt{cabinets.id}, non nul, indexé, suppression en cascade. \\
Données rattachées à un dossier : \texttt{entreprise\_id} & \texttt{UUID} & FK vers \texttt{entreprises.id}, le plus souvent non nulle et indexée ; elle est facultative dans \texttt{documents} et \texttt{taches}. \\
\shortstack[l]{\texttt{chronos.}\\\texttt{entreprise\_id}} & \texttt{UUID} & FK non nulle et \texttt{UNIQUE} : un chrono principal au maximum par entreprise. \\
\shortstack[l]{\texttt{documents.}\\\texttt{entreprise\_id}\\\texttt{documents.}\\\texttt{chrono\_id}} & \texttt{UUID NULL} & Clés étrangères facultatives pour conserver une pièce propre au cabinet. \\
\shortstack[l]{\texttt{documents.}\\\texttt{donnees\_extraites}} & \texttt{JSONB} & Données originales et enrichies du pipeline. \\
\shortstack[l]{\texttt{documents.}\\\texttt{texte\_ocr}} & \texttt{TEXT} & Texte extrait conservé pour contrôle et audit. \\
\shortstack[l]{\texttt{documents.}\\\texttt{hash\_fichier}} & \texttt{VARCHAR(64)} & Index de détection des doublons binaires. \\
Montants extraits de facture & \texttt{NUMERIC(12,2)} & HT, TVA et TTC conservés avec deux décimales. \\
Débit/crédit, allocations et valeurs MAD converties & \texttt{NUMERIC(14,2)} & Précision comptable, valeurs non flottantes. \\
Montants en devise & \texttt{NUMERIC(18,6)} & Montants originaux et parts affectées conservés avec une précision supérieure. \\
Cours de change et cours BAM & \texttt{NUMERIC(20,10)} & Conservation des taux sans troncature prématurée. \\
Statuts et catégories & \texttt{ENUM} ou \texttt{VARCHAR} & Vocabulaires contrôlés par modèle, schéma et service. \\
Dates métier & \texttt{DATE} & Date de pièce, opération, cours ou régularisation. \\
Dates d'audit & \texttt{TIMESTAMP WITH TIME ZONE} & Création, modification, validation et confirmation. \\
\shortstack[l]{\texttt{lignes\_}\\\texttt{comptables}} & trois \texttt{CHECK} & Montants positifs, un seul sens débit/crédit et une seule origine parmi facture, banque ou clôture. \\
Plan comptable & \shortstack[l]{\texttt{UNIQUE(cabinet\_id,}\\\texttt{entreprise\_id,}\\\texttt{numero\_compte)}} & Un numéro est unique uniquement dans le dossier concerné. \\
Allocations bancaires & \shortstack[l]{\texttt{UNIQUE(}\\\texttt{mouvement\_}\\\texttt{bancaire\_id,}\\\texttt{ecriture\_id)}} & Empêche de dupliquer une même paire de rapprochement. \\
Périodes TVA & \shortstack[l]{\texttt{UNIQUE(cabinet\_id,}\\\texttt{entreprise\_id,}\\\texttt{annee, mois)}} & Une période par mois et par dossier. \\
Taux BAM & \shortstack[l]{\texttt{UNIQUE(date\_cours,}\\\texttt{devise, source)}} & Évite la duplication d'un cours historique. \\
\bottomrule
\end{longtable}
\normalsize

\begin{figure}[H]
\centering
\begin{tikzpicture}[node distance=0.45cm]
  \node[layer, fill=blue!6] (mcd) {MCD : entités, associations, propriétés et cardinalités métier};
  \node[layer, fill=green!7, below=of mcd] (mld) {MLD : relations, clés primaires, clés étrangères et unicités};
  \node[layer, fill=ComptaBlue!5, below=of mld] (mpd) {MPD PostgreSQL : UUID, NUMERIC, JSONB, index, CHECK et actions référentielles};
  \node[layer, fill=ComptaGray!6, below=of mpd] (orm) {Implémentation : modèles SQLAlchemy + migrations Alembic};
  \draw[flow] (mcd) -- (mld);
  \draw[flow] (mld) -- (mpd);
  \draw[flow] (mpd) -- (orm);
\end{tikzpicture}
\caption{Passage de la conception Merise à l'implémentation}
\label{fig:merise-passage}
\end{figure}

\subsection{Dictionnaire synthétique des tables}

Afin d'éviter de répéter le MLD complet, ce dictionnaire est volontairement limité aux tables du cœur documentaire, comptable et bancaire représentées dans le MCD. Les tables spécialisées de TVA, de change, de clôture, de tâches et d'audit restent recensées dans le MLD.

\small
\TableSetup
\begin{longtable}{|L{5.4cm}|L{9cm}|}
\caption{Tables principales de ComptaFlow}\label{tab:tables}\\
\toprule
\rowcolor{ComptaBlue}\TableHead{Table} & \TableHead{Responsabilité} \\
\midrule
\endfirsthead
\toprule
\rowcolor{ComptaBlue}\TableHead{Table} & \TableHead{Responsabilité} \\
\midrule
\endhead
\texttt{cabinets} & Tenant principal, coordonnées et état du cabinet. \\
\texttt{users} & Utilisateurs, identité, rôle, état et dernier accès. \\
\texttt{entreprises} & Dossiers comptables gérés par un cabinet. \\
\texttt{chronos} & Regroupement documentaire par entreprise. \\
\texttt{documents} & Métadonnées du fichier, empreinte, statut, OCR, JSON extrait et erreurs. \\
\texttt{ecritures\_comptables} & Pré-écriture comptable, montants, comptes, devise, anomalies et validation. \\
\texttt{lignes\_comptables} & Lignes débit/crédit alimentant Grand Livre, Balance et états. \\
\shortstack[l]{\texttt{comptes\_comptables\_}\\\texttt{entreprise}} & Plan exact d'une entreprise, famille CGNC, usage et nature. \\
\shortstack[l]{\texttt{comptes\_bancaires\_}\\\texttt{entreprise}} & Banque, RIB/IBAN, devise et compte comptable associé. \\
\texttt{mouvements\_bancaires} & Opérations extraites, conversion, compte, rapprochement et virement lié. \\
\shortstack[l]{\texttt{rapprochements\_bancaires\_}\\\texttt{allocations}} & Relation plusieurs-à-plusieurs entre paiements et factures avec montant alloué. \\
\bottomrule
\end{longtable}
\normalsize

\section{Conception applicative}

L'application adopte une architecture modulaire afin de séparer l'interface
utilisateur, les échanges avec l'API, les règles métier, les traitements
asynchrones et la persistance. Cette organisation facilite la maintenance et
isole les traitements documentaires lourds du fonctionnement courant de
l'interface.

\subsection{Architecture en couches}

\begin{figure}[H]
\centering
\begin{tikzpicture}[
  scale=0.82,
  every node/.style={transform shape},
  font=\small,
  appcard/.style={
    rectangle split,
    rectangle split parts=3,
    rounded corners=3pt,
    draw=ComptaBlue!55!ComptaDark,
    line width=0.65pt,
    text width=3.65cm,
    minimum width=4.15cm,
    minimum height=2.75cm,
    align=center,
    inner sep=5pt
  },
  archflow/.style={-{Stealth[length=2.3mm]}, thick, draw=ComptaBlue!75!ComptaDark},
  zone/.style={rounded corners=5pt, draw=ComptaGray!75, dashed, line width=0.55pt, inner sep=7pt}
]
  \node[appcard, fill=ComptaBlue!5] (presentation) at (0,5.15)
    {\textbf{\color{ComptaDark}Présentation}
     \nodepart{two}Interface utilisateur\\Pages, composants et formulaires
     \nodepart{three}\scriptsize\color{ComptaBlue}React -- TypeScript -- Vite\\Tailwind CSS};

  \node[appcard, fill=ComptaBlue!4] (communication) at (4.95,5.15)
    {\textbf{\color{ComptaDark}Communication API}
     \nodepart{two}Requêtes HTTP\\Gestion du jeton et des erreurs
     \nodepart{three}\scriptsize\color{ComptaBlue}Axios};

  \node[appcard, fill=ComptaGreen!7] (api) at (9.9,5.15)
    {\textbf{\color{ComptaDark}API et contrôle d'accès}
     \nodepart{two}Routes REST, authentification\\Rôles et validation des entrées
     \nodepart{three}\scriptsize\color{ComptaGreen!55!black}FastAPI -- JWT / RBAC -- Pydantic};

  \node[appcard, fill=ComptaGreen!6] (services) at (9.9,0)
    {\textbf{\color{ComptaDark}Services métier}
     \nodepart{two}Documents, entreprises, chronos\\Comptabilité, banque, TVA et clôture
     \nodepart{three}\scriptsize\color{ComptaGreen!55!black}Services Python};

  \node[appcard, fill=ComptaGray!6] (async) at (4.95,0)
    {\textbf{\color{ComptaDark}Traitements asynchrones}
     \nodepart{two}Extraction et classification\\OCR / IA et reprises
     \nodepart{three}\scriptsize\color{ComptaSlate}Celery -- Redis};

  \node[appcard, fill=ComptaGray!12] (persistence) at (0,0)
    {\textbf{\color{ComptaDark}Persistance}
     \nodepart{two}Données, historique\\PDF et images stockés localement
     \nodepart{three}\scriptsize\color{ComptaDark}PostgreSQL -- SQLAlchemy\\Alembic -- stockage local};

  \draw[archflow] (presentation.east) -- (communication.west);
  \draw[archflow] (communication.east) -- (api.west);
  \draw[archflow] (api.south) -- (services.north);
  \draw[archflow] (services.west) -- (async.east);
  \draw[archflow] (async.west) -- (persistence.east);

  \node[zone, fit=(presentation)(communication)] (interfacezone) {};
  \node[zone, fit=(api)(services)] (corezone) {};
  \node[zone, fit=(async)(persistence)] (infrazone) {};
  \node[font=\scriptsize\bfseries, text=ComptaGray!45!ComptaDark,
        fill=white, inner sep=1.5pt] at (interfacezone.north) {Interface utilisateur};
  \node[font=\scriptsize\bfseries, text=ComptaGray!45!ComptaDark,
        fill=white, inner sep=1.5pt] at (corezone.north) {Cœur applicatif};
  \node[font=\scriptsize\bfseries, text=ComptaGray!45!ComptaDark,
        fill=white, inner sep=1.5pt] at (infrazone.north) {Infrastructure};
\end{tikzpicture}
\caption{Architecture en couches de l'application}
\label{fig:layers}
\end{figure}

\subsection{Diagramme de composants}

Le diagramme suivant précise les principaux composants techniques et les
échanges entre eux.

\begin{figure}[H]
\centering
\begin{tikzpicture}[
  font=\small,
  component/.style={rounded corners=4pt, draw=ComptaBlue!60!ComptaDark,
    line width=0.7pt, fill=white, align=center, minimum height=0.9cm,
    text width=3.55cm, inner sep=5pt},
  infrastructure/.style={component, draw=ComptaGreen!60!black, fill=ComptaGreen!5,
    text width=3.25cm},
  direct/.style={-{Stealth[length=2.4mm]}, thick, draw=ComptaBlue!75!ComptaDark},
  directboth/.style={{Stealth[length=2.4mm]}-{Stealth[length=2.4mm]}, thick,
    draw=ComptaBlue!75!ComptaDark},
  asynchronous/.style={-{Stealth[length=2.4mm]}, thick, dashed,
    draw=ComptaSlate},
  asynchronousboth/.style={{Stealth[length=2.4mm]}-{Stealth[length=2.4mm]}, thick,
    dashed, draw=ComptaSlate},
  edgelabel/.style={font=\scriptsize, fill=white, inner sep=1.5pt, text=ComptaDark}
]
  \node[draw=none, fill=none, font=\small\bfseries, text=ComptaDark] (user) at (0,13.5)
    {Utilisateur / navigateur};
  \node[component, fill=ComptaBlue!5] (frontend) at (0,11.7)
    {\textbf{Frontend Web}\\React, TypeScript, Vite, Tailwind CSS};
  \node[component, fill=ComptaGreen!6] (fastapi) at (0,9.35)
    {\textbf{API FastAPI}\\Routes REST, JWT / RBAC, Pydantic};
  \node[component, fill=ComptaGreen!5, text width=4.45cm] (services) at (0,6.65)
    {\textbf{Services métier}\\Documents, entreprises, chronos, comptabilité,\\banque, TVA, états et clôture};

  \node[infrastructure] (postgres) at (-4.55,3.25)
    {\textbf{PostgreSQL}\\Données métier et historique};
  \node[infrastructure, fill=ComptaGray!6, draw=ComptaSlate] (redis) at (0,3.25)
    {\textbf{Redis}\\Broker Celery et backend de résultats};
  \node[infrastructure] (storage) at (4.55,3.25)
    {\textbf{Stockage documentaire local}\\PDF et images};
  \node[infrastructure, fill=ComptaGray!6, draw=ComptaSlate] (worker) at (0,-0.1)
    {\textbf{Worker Celery}\\Traitement documentaire asynchrone};
  \node[infrastructure, text width=4.4cm] (ocria) at (4.8,-0.1)
    {\textbf{Extraction / OCR / IA}\\Groq Vision et texte\\Replis : texte PDF natif, PaddleOCR, Regex, Ollama};

  \draw[direct] (user) -- (frontend);
  \draw[directboth] (frontend) -- node[edgelabel, right]{HTTPS / JSON} (fastapi);
  \draw[direct] (fastapi) -- node[edgelabel, right]{appels directs} (services);
  \draw[directboth] (services.south west) -- node[edgelabel, above, sloped]{lecture / écriture} (postgres.north);
  \draw[directboth] (fastapi.east) -| node[edgelabel, near end, right]{fichiers} (storage.north);
  \draw[asynchronous] (fastapi.west) -- ++(-1.35,0) |- node[edgelabel, near end, above]{tâches Celery} (redis.north);
  \draw[asynchronousboth] (redis) -- node[edgelabel, right]{broker / résultats} (worker);
  \draw[directboth] (worker.west) -- node[edgelabel, below, sloped]{lecture / écriture} (postgres.south);
  \draw[direct] (storage.south west) -- node[edgelabel, below, sloped]{lecture} (worker.north east);
  \draw[direct] (worker) -- node[edgelabel, above]{appels} (ocria);

  \draw[direct] (-4.7,-2.05) -- ++(0.75,0)
    node[anchor=west, font=\scriptsize, text=ComptaDark]{appel ou accès direct};
  \draw[asynchronous] (0.2,-2.05) -- ++(0.75,0)
    node[anchor=west, font=\scriptsize, text=ComptaDark]{échange asynchrone};
\end{tikzpicture}
\caption{Diagramme des composants techniques}
\label{fig:components}
\end{figure}

\subsection{Organisation du code}

L'organisation du code reprend la séparation des responsabilités définie par
l'architecture applicative.

\begin{figure}[htbp]
\centering
\begin{tikzpicture}[
  codepanel/.style={rounded corners=5pt, line width=0.8pt, fill=white,
    minimum width=7.15cm, minimum height=10.1cm, inner sep=8pt,
    align=left, anchor=north west}
]
  \node[codepanel, draw=ComptaGreen!65!black] (backend) at (0,0) {%
    \begin{minipage}[t]{6.55cm}
    {\large\bfseries\color{ComptaGreen!55!black}BACKEND}\par
    \vspace{0.25em}
    {\small\color{ComptaGray!45!ComptaDark}\texttt{backend/app/}}\par
    \vspace{0.5em}
    \renewcommand{\arraystretch}{1.42}
    \begin{tabular}{@{}>{\raggedright\arraybackslash}p{2.45cm}>{\raggedright\arraybackslash}p{3.65cm}@{}}
      \texttt{api/}      & Routes HTTP FastAPI \\
      \texttt{core/}     & Configuration, sécurité, base et Celery \\
      \texttt{models/}   & Modèles SQLAlchemy \\
      \texttt{schemas/}  & Schémas Pydantic \\
      \texttt{services/} & Logique métier et intégrations \\
      \texttt{tasks/}    & Traitements asynchrones \\
      \texttt{main.py}   & Initialisation de l'API et des routeurs \\
    \end{tabular}
    \end{minipage}
  };

  \node[codepanel, draw=ComptaBlue!70!ComptaDark] (frontend) at (7.65,0) {%
    \begin{minipage}[t]{6.55cm}
    {\large\bfseries\color{ComptaBlue}FRONTEND}\par
    \vspace{0.25em}
    {\small\color{ComptaGray!45!ComptaDark}\texttt{frontend/src/}}\par
    \vspace{0.5em}
    \renewcommand{\arraystretch}{1.31}
    \begin{tabular}{@{}>{\raggedright\arraybackslash}p{2.45cm}>{\raggedright\arraybackslash}p{3.65cm}@{}}
      \texttt{api/}        & Appels Axios vers le backend \\
      \texttt{assets/}     & Ressources statiques \\
      \texttt{components/} & Composants réutilisables \\
      \texttt{context/}    & État partagé d'authentification \\
      \texttt{pages/}      & Vues principales \\
      \texttt{types/}      & Types TypeScript \\
      \texttt{utils/}      & Fonctions utilitaires \\
      \texttt{App.tsx}     & Routage et composition des pages \\
      \texttt{main.tsx}    & Point d'entrée React \\
    \end{tabular}
    \end{minipage}
  };

  \draw[ComptaGreen!65!black, line width=2.2pt]
    ([xshift=0.16cm,yshift=-0.18cm]backend.north west) --
    ([xshift=0.16cm,yshift=0.18cm]backend.south west);
  \draw[ComptaBlue, line width=2.2pt]
    ([xshift=0.16cm,yshift=-0.18cm]frontend.north west) --
    ([xshift=0.16cm,yshift=0.18cm]frontend.south west);
\end{tikzpicture}
\caption{Organisation du code du projet}
\label{fig:code-organization}
\end{figure}

\subsection{Sécurité et confidentialité}

La sécurité est appliquée à plusieurs niveaux complémentaires. L'API vérifie
l'identité et le rôle de l'utilisateur, tandis que les requêtes métier
appliquent les portées du cabinet et de l'entreprise lorsque le dossier est
concerné. Les contrôles documentaires interviennent avant l'enregistrement du
fichier et les actions sensibles prises en charge par la piste d'audit
conservent leur auteur et les états avant/après disponibles.

\begin{table}[H]
\centering
\caption{Mécanismes de sécurité et de confidentialité appliqués}
\footnotesize
\TableSetup
\begin{tabularx}{\textwidth}{|L{3.8cm}|Y|}
\toprule
\rowcolor{ComptaBlue}\TableHead{Domaine} & \TableHead{Mécanisme appliqué} \\
\midrule
Authentification & Jetons JWT signés en HS256, durée de validité de huit heures et rejet des jetons invalides ou expirés. \\
Mots de passe & Hachage bcrypt et vérification centralisée ; aucun mot de passe en clair n'est utilisé pour l'authentification. \\
Rôles et comptes & RBAC pour les opérations sensibles et refus des comptes marqués inactifs. \\
Isolation multi-tenant & Filtrage par \texttt{cabinet\_id} et contrôle de \texttt{entreprise\_id} pour les données propres aux dossiers comptables. \\
Fichiers importés & PDF, PNG et JPEG uniquement, taille maximale de 20~Mo et cohérence contrôlée entre signature binaire, type MIME déclaré et extension. \\
Intégrité et doublons & Empreinte SHA-256 et recherche d'un fichier identique dans le cabinet avant l'enregistrement. \\
Stockage & Répertoire documentaire séparé par cabinet et refus de résoudre un chemin situé hors de la racine configurée. \\
Traçabilité & Journal d'audit pour les opérations sensibles instrumentées, avec utilisateur, cabinet, ressource et états avant/après ; cette journalisation n'est pas présentée comme exhaustive. \\
Configuration & Clé de signature, accès PostgreSQL, Redis, Ollama et clé Groq fournis par variables d'environnement, sans inscription de secrets dans le code. \\
\bottomrule
\end{tabularx}
\end{table}

L'utilisation d'un service d'IA distant comme Groq implique, dans un contexte
réel de production, de définir une politique de confidentialité adaptée aux
documents comptables traités. Le projet ne revendique ni certification ni
accord contractuel particulier sur ce point. Les captures reproduites dans ce
rapport utilisent exclusivement des données fictives créées pour la
démonstration.

% ============================================================
% 10. Réalisation
% ============================================================
\chapter{Réalisation}

\section{Technologies utilisées}

\begin{table}[H]
\centering
\caption{Technologies du projet}
\footnotesize
\TableSetup
\begin{tabularx}{\textwidth}{|L{3.4cm}|L{4cm}|Y|}
\toprule
\rowcolor{ComptaBlue}\TableHead{Domaine} & \TableHead{Technologie} & \TableHead{Rôle} \\
\midrule
Frontend & React 19, TypeScript, Vite & SPA, typage, compilation et chargement différé des pages. \\
Interface & Tailwind CSS, Lucide React & Mise en forme responsive et iconographie. \\
Communication & Axios & Client HTTP et interception des erreurs d'authentification. \\
Backend & Python (cible Docker 3.12), FastAPI, Pydantic & API REST, validation des entrées et sérialisation. \\
Métier et données & Python, SQLAlchemy & Règles comptables et accès objet-relationnel. \\
Base & PostgreSQL & Persistance transactionnelle multi-tenant. \\
Migrations & Alembic & Évolution versionnée du schéma. \\
Asynchrone & Celery, Redis & File de tâches et exécution du pipeline lourd. \\
OCR et image & PaddleOCR, PaddlePaddle, OpenCV, Pillow & Prétraitement et extraction locale du texte. \\
IA & Groq Vision/texte, Ollama & Extraction/classification principale et repli local. \\
Documents & ReportLab, Pandas, OpenPyXL & Rapports et exports tabulaires. \\
Tests & Pytest & Validation automatisée des services et règles. \\
Déploiement & Docker, Docker Compose, Nginx & Conteneurisation, orchestration et diffusion du frontend. \\
\bottomrule
\end{tabularx}
\end{table}

\section{Architecture du logiciel}

L'architecture définie lors de la phase de conception a été mise en œuvre sous
la forme d'une application web séparant l'interface utilisateur, l'API, les
services métier, les traitements asynchrones et la persistance. Le frontend,
développé avec React et TypeScript puis compilé avec Vite, constitue une
application distincte du backend FastAPI.

Le frontend communique avec le backend par l'intermédiaire d'une API REST et
du client HTTP Axios. Pour les requêtes protégées, le jeton JWT est transmis
dans l'en-tête d'autorisation. FastAPI\cite{fastapi} vérifie alors l'identité de
l'utilisateur, son état et, lorsque l'opération l'exige, son rôle avant de
déléguer le traitement aux services métier. Ces services regroupent notamment
les règles relatives aux documents, aux entreprises, aux chronos, à la
comptabilité et aux opérations bancaires, tout en appliquant la portée du
cabinet concerné.

Après la connexion, le routage frontend conduit d'abord vers une page d'accueil
intermédiaire et non vers un dossier d'entreprise présélectionné. Cette page
sépare explicitement l'import de nouvelles pièces de l'ouverture d'un dossier
existant. Dans le second parcours, l'utilisateur recherche une entreprise par
nom, ICE, IF ou RC, puis sélectionne son contexte actif. Cet identifiant facilite
la navigation entre les écrans, mais ne constitue jamais une autorisation :
chaque appel API contrôle à nouveau le cabinet, l'entreprise et les permissions
de l'utilisateur.

Les traitements documentaires susceptibles de durer davantage, notamment
l'extraction, l'OCR, la classification et les analyses associées, sont confiés
à un worker Celery\cite{celery}. Redis\cite{redis} assure à la fois le transport des tâches et le
stockage de leurs résultats. Cette exécution asynchrone évite de bloquer
l'interface pendant le traitement et permet à l'application de suivre ensuite
l'état du document ainsi que les résultats produits.

La persistance des données métier repose sur PostgreSQL\cite{postgresql}, auquel l'application
accède par SQLAlchemy\cite{sqlalchemy}, tandis qu'Alembic gère les évolutions versionnées du
schéma. Les pièces importées sont conservées dans le stockage documentaire
local configuré par l'application, dans un répertoire propre au cabinet. Cette
séparation des responsabilités facilite la maintenance et permet de faire
évoluer les différents composants de manière indépendante.

\section{Pipeline complet de traitement}

Le traitement d'une pièce comptable repose sur une chaîne successive de
contrôles, d'extraction, d'identification et d'application des règles métier.
Lors de l'import, l'API contrôle le contenu du fichier, calcule son empreinte,
écarte les doublons binaires déjà connus dans le cabinet, puis enregistre le
document dans le stockage local. Une tâche Celery, transmise par Redis, prend
ensuite en charge les opérations documentaires sans bloquer l'interface.

Le système tente d'abord une extraction et une classification unifiées par
Groq Vision\cite{groq}. Si ce traitement n'aboutit pas, il exploite la couche texte native
du PDF ou, à défaut, PaddleOCR\cite{paddleocr}, avant d'appliquer les mécanismes d'analyse
textuelle prévus. Le repli local peut également mobiliser Ollama\cite{ollama}. Les valeurs obtenues sont normalisées et contrôlées sans
remplacer silencieusement les montants extraits.

L'entreprise concernée, le tiers et le sens achat/vente sont ensuite déterminés
à partir des acteurs présents sur la pièce et des dossiers réellement suivis
par le cabinet. Après la détection des doublons métier et l'affectation au
chrono, une facture conduit à la préparation d'une pré-écriture comptable,
tandis qu'un relevé bancaire conduit à la création de mouvements pouvant être
rapprochés avec les factures correspondantes. Les résultats incomplets ou
ambigus restent soumis à un contrôle humain.

Les lignes comptables sont préparées dès le stade de brouillon pour permettre
leur prévisualisation. Le cycle distingue ensuite le calcul en cours, les
éléments à vérifier, les pré-écritures prêtes pour Topaze et celles dont la
saisie dans Topaze a été confirmée manuellement. Seules les données ayant
franchi les contrôles et la validation humaine alimentent les journaux, le
Grand Livre, la Balance, la TVA, le CPC et le Bilan.

\begin{figure}[htbp]
\centering
\begin{tikzpicture}[
  scale=0.85,
  every node/.style={transform shape},
  font=\small,
  pipebox/.style={rounded corners=4pt, draw=ComptaBlue!65!ComptaDark,
    line width=0.7pt, fill=white, align=center, text width=3.2cm,
    minimum height=1.45cm, inner sep=4pt},
  analysisbox/.style={pipebox, fill=ComptaBlue!5},
  businessbox/.style={pipebox, fill=ComptaGreen!6,
    draw=ComptaGreen!65!black},
  decision/.style={diamond, aspect=2.25, draw=ComptaSlate,
    line width=0.75pt, fill=ComptaGray!6, align=center,
    text width=2.3cm, inner sep=2pt},
  pipeflow/.style={-{Stealth[length=2.4mm]}, thick,
    draw=ComptaBlue!75!ComptaDark},
  phase/.style={rounded corners=6pt, draw=ComptaGray!80, dashed,
    line width=0.55pt, inner sep=8pt},
  phaselabel/.style={font=\small\bfseries,
    text=ComptaGray!40!ComptaDark, fill=white, inner sep=2pt},
  edgeword/.style={font=\footnotesize\bfseries,
    text=ComptaDark, fill=white, inner sep=1pt}
]
  % Phase A : acquisition documentaire
  \node[pipebox, fill=ComptaLight] (upload) at (-5.7,15.3)
    {\textbf{1. Import du\\document}};
  \node[pipebox] (control) at (-1.9,15.3)
    {\textbf{2. Contrôle du\\fichier}\\[-0.1em]
     \footnotesize SHA-256, doublon exact, taille 20 Mo\\MIME, signature et extension};
  \node[pipebox] (storage) at (1.9,15.3)
    {\textbf{3. Stockage du\\document}\\[-0.1em]
     \footnotesize stockage local par cabinet\\statut \texttt{en\_attente}};
  \node[pipebox, fill=ComptaGray!6, draw=ComptaSlate] (queue) at (5.7,15.3)
    {\textbf{4. Mise en file\\du traitement}\\[-0.1em]
     \footnotesize Celery / Redis};

  % Phase B : extraction et analyse
  \node[analysisbox, text width=4.7cm] (extract) at (0,12.45)
    {\textbf{5. Extraction et\\classification principales}\\[-0.1em]
     \footnotesize Groq Vision};
  \node[decision] (usable) at (0,10.35)
    {{\NoAutoSpacing\textbf{Données\\exploitables ?}}};
  \node[analysisbox, fill=ComptaGray!6, draw=ComptaSlate,
        text width=4.45cm] (fallback) at (-3.9,8.1)
    {\textbf{6. Mécanismes\\de repli}\\[-0.1em]
     \footnotesize texte PDF natif\\ou PaddleOCR\\Groq texte / Regex / ML\\puis Ollama};
  \node[analysisbox, text width=4.05cm] (kept) at (3.9,8.1)
    {\textbf{Conservation des\\données extraites}\\[-0.1em]
     \footnotesize valeurs et source d'extraction};
  \node[analysisbox, text width=4.8cm] (normalize) at (0,5.85)
    {\textbf{7. Normalisation et\\contrôle de cohérence}\\[-0.1em]
     \footnotesize HT / TVA / TTC ou totaux et solde bancaire};
  \node[analysisbox, text width=4.5cm] (identify) at (0,3.55)
    {\textbf{8. Identification\\métier}\\[-0.1em]
     \footnotesize entreprise, tiers et sens achat/vente};

  % Phase C : traitement métier et validation
  \node[businessbox] (chrono) at (-5.7,0.25)
    {\textbf{9. Doublon\\métier et chrono}\\[-0.1em]
     \footnotesize avant toute pré-écriture};
  \node[businessbox] (rules) at (-1.9,0.25)
    {\textbf{10. Règles métier\\et comptes}\\[-0.1em]
     \footnotesize plan comptable de l'entreprise};
  \node[businessbox] (prepare) at (1.9,0.25)
    {\textbf{11. Préparation\\comptable}\\[-0.1em]
     \footnotesize pré-écriture ou traitement bancaire};
  \node[businessbox] (human) at (5.7,0.25)
    {\textbf{12. Contrôle\\humain}\\[-0.1em]
     \footnotesize correction,\\validation ou rejet};
  \node[businessbox, fill=ComptaLight, text width=5.8cm] (statements) at (0,-2.55)
    {\textbf{13. Exploitation des\\données validées}\\[-0.1em]
     \footnotesize Journaux, Grand Livre,\\Balance, TVA, CPC et Bilan};

  % Flux principal et branche de repli
  \draw[pipeflow] (upload) -- (control);
  \draw[pipeflow] (control) -- (storage);
  \draw[pipeflow] (storage) -- (queue);
  \draw[pipeflow] (queue.south) |- (extract.east);
  \draw[pipeflow] (extract) -- (usable);
  \draw[pipeflow] (usable.west) -| node[edgeword, pos=0.25, above]{non} (fallback.north);
  \draw[pipeflow] (usable.east) -| node[edgeword, pos=0.25, above]{oui} (kept.north);
  \draw[pipeflow] (fallback.south) |- (normalize.west);
  \draw[pipeflow] (kept.south) |- (normalize.east);
  \draw[pipeflow] (normalize) -- (identify);
  \draw[pipeflow] (identify.south) -- ++(0,-0.4) -| (chrono.north);
  \draw[pipeflow] (chrono) -- (rules);
  \draw[pipeflow] (rules) -- (prepare);
  \draw[pipeflow] (prepare) -- (human);
  \draw[pipeflow] (human.south) |- (statements.east);

  % Regroupements visuels
  \node[phase, fit=(upload)(control)(storage)(queue)] (phaseA) {};
  \node[phase, fit=(extract)(usable)(fallback)(kept)(normalize)(identify)] (phaseB) {};
  \node[phase, fit=(chrono)(rules)(prepare)(human)(statements)] (phaseC) {};
  \node[phaselabel] at (phaseA.north) {A. Acquisition documentaire};
  \node[phaselabel] at (phaseB.north) {B. Extraction et analyse};
  \node[phaselabel] at (phaseC.north) {C. Traitement métier et validation};
\end{tikzpicture}
\caption{Pipeline complet de traitement d'une pièce comptable}
\label{fig:pipeline}
\end{figure}

Le pipeline suit ainsi une progression contrôlée : les informations extraites
ne sont pas considérées comme définitives. Elles passent par des contrôles de
cohérence, une identification métier et, lorsque la situation le nécessite,
une intervention humaine. Cette organisation associe automatisation du
traitement documentaire et contrôle professionnel des résultats.

\subsection{Workflow pré-comptable et passage vers Topaze}

Le service de workflow exécute des contrôles déterministes sur les valeurs
obligatoires, la date, la cohérence HT--TVA--TTC, l'équilibre Débit/Crédit, les
doublons et l'existence de chaque compte exact dans le plan de l'entreprise.
Une pré-écriture conforme passe au statut \texttt{prete\_topaze}. Une ambiguïté
ou un compte absent la maintient au statut \texttt{a\_verifier}, avec une
anomalie structurée et sans ligne fictive d'équilibrage. La saisie réalisée
dans Topaze reste une action externe : l'utilisateur autorisé la confirme dans
ComptaFlow, qui conserve la date, l'auteur et une référence de lot facultative.

La TVA suit la même séparation. La synthèse mensuelle est calculée depuis les
lignes validées du Grand Livre. La centralisation n'est préparée que si les
comptes de TVA exacts ont été configurés et existent dans le plan de
l'entreprise ; elle doit être strictement équilibrée. Les régularisations ou
retenues dont la contrepartie n'est pas déterminable restent à vérifier. Les
statuts comptable, déclaratif et Topaze sont conservés séparément afin de ne
pas confondre calcul interne, déclaration fiscale et saisie dans le logiciel de
référence.

Enfin, une période de travail peut être verrouillée après les contrôles de
pré-clôture. Elle devient alors accessible en lecture seule : les mutations
comptables, bancaires, fiscales et documentaires portant sur cette période sont
refusées avant toute modification ou mise en file d'un traitement. Les
consultations restent disponibles. Une réouverture nécessite un rôle autorisé,
une justification et une trace d'audit avant que les mutations redeviennent
possibles.

\subsection{Traitement des factures}

Lors du traitement d'une facture, les valeurs extraites sont conservées dans
les données du document afin de préserver la source du traitement. Le service
d'identification recherche séparément le fournisseur et le client parmi les
dossiers confirmés du cabinet. La correspondance suit l'ordre ICE, identifiant
fiscal, RC puis nom normalisé. Les champs historiques
\texttt{nom\_entreprise} et \texttt{tiers} ne sont utilisés comme émetteur et
destinataire que lorsque les deux rôles structurés sont absents.

Le sens de l'opération est déterminé depuis le point de vue du dossier suivi.
Si l'entreprise gérée est cliente ou destinataire, la facture est classée en
achat et le tiers correspond au fournisseur. Si elle est fournisseur ou
émettrice, la facture est classée en vente et le tiers correspond au client.
Lorsque les deux parties correspondent à des dossiers suivis, ou qu'aucun
dossier ne peut être reconnu de façon sûre, la direction reste à vérifier au
lieu d'être choisie arbitrairement.

Après le contrôle du doublon métier et l'affectation au chrono, le service crée
ou met à jour la pré-écriture comptable associée au document. Les lignes
Débit/Crédit sont préparées pour la prévisualisation, mais leur statut reste
\texttt{brouillon} ou \texttt{a\_verifier}. Une correction replace la
pré-écriture à vérifier, et la validation est refusée tant que ses lignes ne
sont pas complètes et équilibrées.

\subsection{Résolution des comptes}

Les moteurs d'extraction ne doivent pas fournir un numéro de compte définitif :
ils proposent une nature économique appartenant à un vocabulaire contrôlé.
Les règles Python transforment cette nature en famille CGNC\cite{cgnc}, puis interrogent
le plan comptable de l'entreprise avec les portées \texttt{cabinet\_id} et
\texttt{entreprise\_id}. La recherche privilégie un compte actif compatible
avec l'usage et, lorsqu'elle est renseignée, avec la nature de l'opération.

Une correction déjà présente sur la pré-écriture reste prioritaire. Le moteur
recherche ensuite un compte unique dans le plan, puis peut accepter une donnée
explicitement fournie si elle appartient à la famille attendue, ou réutiliser
un compte historique confirmé pour la même entreprise et la même nature. Si
plusieurs candidats demeurent, le choix automatique est refusé et leurs
numéros sont conservés pour vérification.

Aucun compte HT ou TVA exact n'est fabriqué lorsque le plan ne permet pas de le
déterminer. Pour les comptes tiers uniquement, un compte ``divers'' configuré
dans le plan, ou un compte provisoire compatible, peut servir de repli ; il
n'est pas assimilé à un compte personnel confirmé du fournisseur ou du client.
Toute absence ou ambiguïté empêchant la construction de lignes complètes place
la pré-écriture au statut \texttt{a\_verifier}.

\subsection{Traitement bancaire}

Les relevés bancaires suivent une branche spécialisée. Chaque ligne extraite est
normalisée puis conservée individuellement, avec son libellé original et son
texte brut dans les données du document. Aucune déduplication n'est appliquée
aux opérations : deux lignes identiques sur le relevé restent deux mouvements
distincts. Les totaux débit/crédit sont comparés aux totaux imprimés et, lorsque
les informations existent, le solde final est rapproché du solde initial et des
mouvements extraits.

L'identification du compte bancaire intervient lors du rapprochement. Le
service recherche d'abord une correspondance non ambiguë par RIB ou IBAN, puis
par nom de banque ; un compte actif unique peut également être retenu. À défaut,
un compte comptable unique portant l'usage \texttt{banque} dans le plan de
l'entreprise peut être utilisé, sans inventer de numéro.

Le rapprochement calcule le montant déjà réglé et le solde restant de chaque
facture. Une correspondance exacte, unique et suffisamment bien notée peut être
affectée automatiquement. Les paiements partiels et les règlements groupés sont
proposés, tandis que les résultats multiples ou ambigus nécessitent une
affectation humaine. La table d'allocations permet ainsi qu'un mouvement règle
plusieurs factures et qu'une facture reçoive plusieurs règlements. Lors d'une
confirmation, la valeur imputée ne peut pas dépasser le solde restant et le
total affecté doit couvrir exactement le mouvement bancaire.

Les virements internes sont liés par paires débit/crédit sur deux comptes
bancaires distincts, avec contrôle du montant et de la proximité des dates.
Les frais bancaires, acomptes et autres opérations spéciales peuvent être
enregistrés avec un compte de contrepartie qui doit exister dans le plan de
l'entreprise ; tant qu'il manque, aucune ligne comptable complète n'est
produite.

\subsection{Gestion des devises et cours BAM}

Pour une opération libellée dans une devise étrangère, le moteur conserve
séparément la devise et les montants d'origine, le montant converti en dirhams,
le taux appliqué, sa date, son type, sa source et l'unité de cotation. Les cours
Bank Al-Maghrib\cite{bam} déjà obtenus sont mis en cache dans PostgreSQL par date, devise
et source, avec leurs valeurs d'achat et de vente.

Pour une facture d'achat, le cours de vente à la clientèle de la date de la
pièce est utilisé ; pour une facture de vente, le système utilise le cours
d'achat auprès de la clientèle. Un débit bancaire suit le cours de vente et un
crédit le cours d'achat. La conversion multiplie le montant par le taux puis
divise le résultat par l'unité de cotation, ce qui prend notamment en compte
les devises publiées par unité de 100.

Les montants extraits ne sont jamais écrasés par leur conversion. Pour un
mouvement bancaire, le montant MAD réellement porté par la banque peut en outre
être conservé séparément de la valeur BAM théorique. Si la date, le sens ou le
cours demandé ne peut pas être établi de façon sûre, le système signale
l'anomalie et refuse toute substitution silencieuse par le cours du jour ; la
pré-écriture ou le mouvement reste alors à vérifier.

\subsection{Exports et interopérabilité}

L'API permet d'exporter les registres validés d'une entreprise, filtrés par
catégorie, année et mois, aux formats CSV, XLSX et PDF. Les fichiers sont
générés en mémoire à partir des mêmes données comptables et de leurs totaux,
puis retournés au client sans stockage permanent d'une copie d'export. Ces
formats facilitent la consultation, le contrôle et les échanges génériques,
mais ne constituent pas un connecteur vers un logiciel tiers.

Aucune intégration automatique directe avec Topaze n'a été réalisée pendant le
stage. Une évolution future pourrait consister à définir un format d'export
compatible ou un mécanisme d'intégration, après validation des contraintes du
logiciel utilisé par SEGURIBAT.

\subsection{Idempotence et reprise des traitements}

L'implémentation limite l'accumulation de résultats lors d'un retraitement
séquentiel. L'import calcule une empreinte SHA-256 et refuse, dans un même
cabinet, un fichier identique déjà enregistré. Pour une facture, le service
recherche la pré-écriture liée au document, réutilise la première occurrence et
supprime d'éventuels doublons historiques avant de synchroniser ses lignes. Pour
un relevé, les mouvements précédemment issus du document sont supprimés puis
recréés dans la même transaction ; un second traitement successif ne les ajoute
donc pas à la série précédente.

La tâche Celery réessaie au maximum trois fois les erreurs considérées comme
transitoires, notamment les délais d'attente, erreurs de connexion, réponses
HTTP~429 et erreurs serveur. Le backoff, la reprise au démarrage du broker et
l'acquittement tardif améliorent la continuité du traitement. En revanche,
l'idempotence stricte n'est pas garantie face à deux workers concurrents : ni
l'empreinte du document ni \texttt{document\_id} dans la table des
pré-écritures ne disposent actuellement d'une contrainte \texttt{UNIQUE}. La
prévention repose donc en partie sur la logique applicative. Un verrou de
traitement ou des contraintes de base adaptées constituent une amélioration à
prévoir avant la production. Le retraitement volontaire réinitialise en outre
les données comptables associées au document ; les identifiants de mouvements
et leurs allocations peuvent alors être recréés.

% ============================================================
% 11. Fonctionnalités principales
% ============================================================
\chapter{Fonctionnalités principales réalisées}

Ce chapitre présente les principales fonctionnalités mises en œuvre dans
l'application. Les interfaces retenues illustrent le parcours documentaire, le
contrôle comptable, le traitement bancaire ainsi que la consultation des
principaux états. Les exemples affichés reposent exclusivement sur des données
fictives créées pour la démonstration.

\section{Authentification et choix du parcours}

L'accès à ComptaFlow commence par une authentification protégée. Après une
connexion réussie, l'utilisateur arrive sur une page principale volontairement
simple. Elle ne surcharge pas l'accueil avec des statistiques : elle lui propose
immédiatement deux parcours distincts, tout en conservant un accès discret aux
notifications et à l'Assistant.

\capture{captures/login-page-verification.png}{Écran d'authentification sécurisé de ComptaFlow}{fig:authentification}

Le premier choix, \og Importer de nouveaux documents \fg{}, conduit directement
à l'espace de dépôt. Le second, \og Travailler sur une entreprise \fg{}, ouvre
la recherche des dossiers disponibles. Un texte d'aide apparaît au survol ou au
focus clavier afin d'expliquer ce second parcours sans alourdir l'écran.

\capture{captures/accueil_focus_entreprise.png}{Accueil à deux parcours et aide contextuelle sur le choix d'une entreprise}{fig:accueil-parcours}

La séparation rend explicites les deux chaînes de navigation suivantes :
\begin{itemize}[leftmargin=1.3cm]
  \item connexion $\rightarrow$ import $\rightarrow$ OCR/IA $\rightarrow$
    extraction $\rightarrow$ identification $\rightarrow$ vérification ou
    correction $\rightarrow$ classement et préparation comptable ;
  \item connexion $\rightarrow$ recherche d'une entreprise $\rightarrow$
    sélection du dossier $\rightarrow$ accès à ses modules métier.
\end{itemize}

La page de sélection filtre dynamiquement les entreprises actives et confirmées
par nom, ICE, IF ou RC. Une fois le dossier choisi, le contexte d'entreprise est
conservé pour la navigation ; les contrôles multi-tenant du backend restent
systématiquement appliqués.

\capture{captures/selection_entreprise.png}{Recherche et sélection du dossier d'une entreprise}{fig:selection-entreprise}

\section{Espace entreprise et tableau de bord}

Une fois l'entreprise sélectionnée, l'utilisateur accède aux modules réellement
disponibles pour ce dossier : tableau de bord, documents, chronos, achats,
ventes, banque, comptes bancaires, écritures, registres, Grand Livre, Balance,
TVA, CPC, Bilan, clôture et pré-clôture. Le tableau de bord centralise les
indicateurs du contexte actif : état des documents, pré-écritures nécessitant
une vérification, situation de la TVA et actions requises.

\capture{captures/dashboard.png}{Tableau de bord du dossier comptable sélectionné}{}

\section{Import et suivi documentaire}

L'interface documentaire permet d'importer une ou plusieurs pièces aux formats
PDF, PNG ou JPEG, puis d'en suivre le statut de traitement. La liste conserve un
accès direct au fichier et à l'écran de vérification correspondant.

\capture{captures/import.png}{Interface d'import d'une pièce comptable}{}

Un document dont l'entreprise n'a pas été reconnue n'est ni supprimé ni masqué.
Il reste conservé au niveau du cabinet et apparaît dans la vue \og Documents à
identifier \fg{}. Un utilisateur autorisé peut consulter les données extraites,
attribuer manuellement une entreprise active du même cabinet, puis déclencher le
retraitement. Cette attribution contrôle également le verrouillage de la période
et ne contourne pas les règles d'isolation.

Le chrono numérique fournit une vue structurée des documents. L'entreprise,
l'année, le mois et la catégorie constituent ici des critères de consultation
portés par les pièces ; ils ne correspondent pas à la création d'un chrono
physique distinct pour chaque période. La recherche et les statuts facilitent le
repérage des éléments à contrôler.

\capture{captures/chronos.png}{Consultation et suivi des documents dans les chronos}{}

\section{Contrôle détaillé et validation}

La vue détaillée rapproche le document original des informations extraites :
entreprise, tiers, numéro, date et montants HT, TVA et TTC. Elle signale les
anomalies sans modifier silencieusement les valeurs ambiguës. L'utilisateur peut
corriger les données et examiner la pré-écriture comptable proposée avant de la
valider ou de la rejeter, selon ses droits.

\capture{captures/document_detail.png}{Contrôle d'un document et de sa pré-écriture comptable}{}

\section{Plan comptable et comptes bancaires}

La préparation comptable s'appuie sur les comptes exacts configurés dans le plan
comptable propre à l'entreprise. La configuration bancaire complète ce mécanisme
en associant à chaque compte une banque, une devise, un RIB, un IBAN et un compte
comptable exact déjà présent dans ce plan. Cette liaison facilite l'identification
du compte concerné lors du traitement des relevés.

\capture{captures/comptes_bancaires.png}{Configuration des comptes bancaires d'une entreprise}{}

\section{Gestion bancaire et rapprochements}

La page Banque présente les mouvements issus des relevés, leurs montants, leur
sens et leur état de rapprochement. Le module prend en charge les rapprochements
simples ainsi que les paiements partiels et les allocations multiples. Une
proposition reste distincte d'un rapprochement confirmé : la validation par un
utilisateur autorisé demeure nécessaire, et le montant affecté ne peut pas
dépasser le solde restant dû.

\capture{captures/banque_v2.png}{Rapprochement des mouvements bancaires}{}

\section{Suivi opérationnel}

Les tâches et rappels complètent le suivi opérationnel par l'affectation d'un
responsable, d'une échéance, d'une priorité et, si nécessaire, d'une récurrence.
L'heure d'échéance est conservée, et la fin d'une tâche récurrente crée la
prochaine occurrence. La vue en tableau répartit les actions entre \og À faire
\fg{}, \og En cours \fg{} et \og Terminée \fg{}, tout en mettant en évidence
les retards. Les relances documentaires et les pré-écritures prêtes mais non
saisies dans Topaze sont réunies dans le même espace.

\capture{captures/taches_planification.png}{Planification et suivi visuel des tâches}{fig:taches-planification}

La messagerie interne permet aux utilisateurs d'échanger avec
l'administrateur du cabinet. Un message reçu peut être transformé en tâche
planifiée avec une date et une heure, sans sortir de la portée du cabinet. Les
notifications regroupent les actions en attente, tandis que les rapports
fournissent une synthèse exportable.

\section{États et contrôles comptables}

Seules les lignes comptables validées alimentent le Grand Livre et la Balance.
Le Grand Livre regroupe les mouvements par compte et permet de contrôler
l'équilibre global ainsi que l'évolution des soldes.

\capture{captures/grand_livre.png}{Consultation du Grand Livre}{}

Le CPC présente les produits, les charges et les résultats calculés pour
l'exercice à partir des données validées. Comme le Bilan et les autres états de
synthèse, il restitue les informations disponibles sans corriger automatiquement
une incohérence comptable.

\capture{captures/cpc_bilan.png}{Consultation des états comptables}{}

\subsection{TVA comptable et préparation Topaze}

La page TVA présente les montants issus des lignes validées, les crédits et les
anomalies de chaque période. Elle permet de configurer les comptes exacts de
centralisation, d'enregistrer l'échéance et la déclaration, puis de confirmer
séparément la saisie effectuée dans Topaze. L'interface rappelle explicitement
qu'il s'agit d'une synthèse comptable préparatoire et non d'une déclaration
fiscale automatique.

\capture{captures/tva_preparation_topaze.png}{TVA comptable et préparation de la centralisation}{fig:tva-topaze}

\subsection{Contrôles de pré-clôture}

La pré-clôture rassemble les contrôles des documents, des pré-écritures, de la
banque, de la TVA et des états comptables. Le score constitue un indicateur
technique ; les anomalies restent classées par niveau et doivent être instruites
avant la clôture.

\capture{captures/precloture.png}{Contrôles de pré-clôture}{}

\section{Assistant sécurisé et journal d'audit}

Dans l'espace authentifié, l'Assistant ComptaFlow est accessible par un bouton
flottant compact. Son activation ouvre un panneau latéral ; une page complète
reste également disponible pour les recherches longues. L'Assistant fournit une recherche en lecture seule sur les
entreprises, documents, pré-écritures, mouvements, TVA et tâches auxquels
l'utilisateur a réellement accès. Le moteur transforme la question en un plan
de requête contrôlé ; il n'exécute pas de SQL fourni par l'utilisateur. Chaque
requête applique \texttt{cabinet\_id} et, lorsqu'un contexte est sélectionné,
\texttt{entreprise\_id}. Les résultats documentaires conservent un lien
sécurisé vers la pièce originale.

\capture{captures/assistant_comptaflow.png}{Assistant ComptaFlow en contexte multi-entreprise}{fig:assistant-comptaflow}

Le journal d'audit est réservé aux administrateurs autorisés. Il centralise les
connexions, validations, changements de statut, actions Topaze, modifications
de configuration TVA, tâches, messages et consultations de l'Assistant. Les
filtres par acteur, entreprise, module, action, statut et période facilitent la
reconstitution d'un traitement, sans exposer les mots de passe, jetons ou
contenus OCR complets.

\capture{captures/historique_audit.png}{Journal d'audit filtrable réservé aux administrateurs}{fig:historique-audit}

\section{Contribution personnelle}

Ce projet a été réalisé individuellement. J'ai pris en charge l'ensemble des
étapes, depuis l'analyse du besoin et la conception fonctionnelle,
informationnelle et applicative jusqu'au développement, à l'intégration, aux
tests et à la validation de l'application.

Ma contribution a couvert à la fois le frontend, développé avec React,
TypeScript et Vite, et le backend fondé sur FastAPI. J'ai réalisé les interfaces
utilisateur, leur communication avec l'API ainsi que les services métier
nécessaires au traitement des documents comptables. Le travail a notamment
porté sur l'import et le suivi documentaire, les chronos, l'extraction et
l'analyse des pièces, l'identification des entreprises et des tiers, la
classification achat/vente et la préparation des pré-écritures comptables avec
maintien d'une validation humaine.

J'ai également développé la gestion du plan comptable propre à chaque
entreprise, les comptes et mouvements bancaires, les rapprochements simples,
partiels ou multiples, ainsi que la consultation du Grand Livre, de la Balance,
de la TVA, du CPC, du Bilan et des contrôles de pré-clôture. La modélisation et
la persistance reposent sur SQLAlchemy et PostgreSQL, avec Alembic pour le suivi
des évolutions du schéma de données.

Enfin, j'ai assuré l'intégration des traitements asynchrones avec Celery et
Redis ainsi que des mécanismes d'extraction OCR et d'analyse assistée par Groq,
PaddleOCR et Ollama. Le travail a compris l'intégration de ces composants, la
résolution des anomalies rencontrées au cours du développement, les tests du
backend et du frontend et la vérification du fonctionnement global de
l'application.

% ============================================================
% 12. Tests et validation
% ============================================================
\chapter{Tests et validation}

\section{Stratégie de validation}

La validation combine des tests automatisés, des contrôles de compilation du frontend et des scénarios fonctionnels complémentaires exécutés soit via l'interface, soit directement au niveau des services métier selon la nature du comportement à vérifier. Les vérifications ciblent surtout les zones à risque : intégrité des montants, équilibre débit/crédit, isolation multi-tenant, résolution des comptes, changements de devise, rapprochements et cohérence des états.

\section{Couverture fonctionnelle des tests automatisés}

La suite Pytest présente dans le dépôt couvre notamment :

\begin{itemize}[leftmargin=*]
  \item les règles de comptes d'achat et de vente et le refus des comptes inconnus ;
  \item la conservation de HT, TVA et TTC et le signalement des incohérences ;
  \item l'identification entreprise/tiers et l'isolation par cabinet ;
  \item la sécurité des fichiers importés et la politique de reprise ;
  \item les écritures équilibrées, Grand Livre et Balance ;
  \item les calculs CPC, Bilan, TVA V2 et pré-clôture ;
  \item les cours BAM, conversions et écarts de change ;
  \item la gestion bancaire avancée, les paiements partiels, les groupes, les allocations et les virements internes ;
  \item le workflow pré-Topaze, la centralisation TVA et le verrouillage des périodes ;
  \item la messagerie, la planification des tâches et leur isolation par cabinet ;
  \item l'Assistant, ses filtres documentaires, ses liens sécurisés et le refus des requêtes SQL ;
  \item la journalisation des actions sensibles et son contrôle d'accès administrateur.
\end{itemize}

\section{Résultats obtenus}

\begin{table}[H]
\centering
\caption{Résultats de validation de la version documentée}
\TableSetup
\begin{tabularx}{\textwidth}{|L{4.6cm}|L{3.4cm}|Y|}
\toprule
\rowcolor{ComptaBlue}\TableHead{Validation} & \TableHead{Résultat} & \TableHead{Observation} \\
\midrule
Import des modèles SQLAlchemy & \TableSuccess{Réussi} & Les modèles du package \texttt{app.models} sont importables. \\
Chaîne Alembic & \TableSuccess{Réussi} & Une seule tête détectée : \texttt{f6d2b0e9a381} ; aucune opération de mise à niveau manquante. \\
Tests backend Pytest & \TableSuccess{Réussi} & 401 tests réussis lors de la validation finale ; aucun échec. \\
Typecheck TypeScript & \TableSuccess{Réussi} & \texttt{npm run typecheck} terminé sans erreur. \\
Tests frontend & \TableSuccess{Réussi} & 18 tests Vitest réussis dans six fichiers ; aucun échec. \\
Build frontend & \TableSuccess{Réussi} & Build Vite de production terminé ; 1\,905 modules transformés. \\
Lint frontend & \textcolor{ComptaSlate}{\textbf{Réussi avec avertissement}} & Aucun échec ; un avertissement Fast Refresh dans \texttt{AuthContext.tsx}. \\
Audit visuel des routes & \TableSuccess{Réussi} & 25 routes authentifiées chargées avec le backend réel ; aucune page vide ni erreur d'accès au backend après correction. \\
Scénarios fonctionnels & \TableSuccess{Réussi} & Dix scénarios métier et de sécurité conformes au comportement attendu. \\
\bottomrule
\end{tabularx}
\end{table}

\section{Exemples de scénarios de recette}

Les scénarios suivants ont été vérifiés à partir du jeu de démonstration et
d'exécutions isolées des services métier. Les modifications temporaires
nécessaires au rapprochement partiel ont été annulées après observation afin de
préserver l'état initial de la base de données.

\small
\TableSetup
\begin{longtable}{|L{0.75cm}|L{3.3cm}|L{4.5cm}|L{5.6cm}|}
\caption{Résultats des scénarios fonctionnels}\label{tab:recette}\\
\toprule
\rowcolor{ComptaBlue}\TableHead{ID} & \TableHead{Scénario} & \TableHead{Résultat attendu} & \TableHead{Résultat réel} \\
\midrule
\endfirsthead
\toprule
\rowcolor{ComptaBlue}\TableHead{ID} & \TableHead{Scénario} & \TableHead{Résultat attendu} & \TableHead{Résultat réel} \\
\midrule
\endhead
T01 & Vérifier le traitement d'une facture d'achat valide déjà importée. & Document traité, montants d'origine conservés et pré-écriture équilibrée préparée. & La facture fictive de Société Démo SARL est enregistrée dans le chrono principal, en achats (08/2026). Les montants HT 5\,000, TVA 1\,000 et TTC 6\,000 MAD sont conservés. La pré-écriture d'achat comporte trois lignes équilibrées, pour un total débit et crédit de 6\,000 MAD, puis passe automatiquement au statut \texttt{PRETE\_TOPAZE} puisque les contrôles ne détectent aucune anomalie. \\
T02 & Importer une facture avec HT + TVA différent de TTC. & Incohérence signalée, valeurs d'origine conservées et aucune correction silencieuse. & Avec HT 1\,000, TVA 200 et TTC 1\,250 MAD, les trois valeurs sont restées inchangées. Le contrôle a produit \og Incohérence entre les montants HT, TVA et TTC extraits \fg{}, a activé \texttt{a\_verifier} et a refusé la génération de lignes déséquilibrées. \\
T03 & Ouvrir une donnée d'un autre cabinet. & La donnée d'un autre cabinet n'est pas accessible. & L'accès à un document de démonstration avec l'identifiant d'un autre cabinet fictif a retourné HTTP 404, \og Document introuvable \fg{}. Aucune ressource n'a été retournée. \\
T04 & Rapprocher un règlement partiel. & Le montant réglé est alloué sans dépasser le solde et le montant restant est conservé. & Dans une transaction de test, une facture de 4\,800 MAD a reçu une allocation de 2\,400 MAD. Le rapprochement a été confirmé en mode partiel et le solde observé était de 2\,400 MAD. La transaction a ensuite été annulée pour restaurer les données de démonstration. \\
T05 & Demander un cours historique introuvable. & L'absence de cours historique est signalée sans substitution silencieuse par un cours courant. & Pour une facture fictive en EUR au 01/01/1900 et une absence contrôlée de cours, le service a retourné le statut \texttt{cours\_introuvable}. Les montants d'origine 1\,000, 200 et 1\,200 sont restés conservés ; les montants comptables MAD sont restés non renseignés et aucun cours courant n'a été substitué. \\
T06 & Générer le Bilan avec un déséquilibre. & Le déséquilibre est signalé sans création de ligne comptable fictive. & Avec un actif de 100 MAD et un passif de 80 MAD, le Bilan a retourné un écart de 20 MAD, \texttt{equilibre = false} et un état à vérifier. Les deux comptes fournis sont restés les seuls comptes du résultat : aucune ligne d'équilibrage n'a été créée. \\
T07 & Préparer puis centraliser une période de TVA. & Seuls les comptes configurés sont utilisés et la centralisation reste équilibrée. & Les tests ont produit les lignes à partir des comptes exacts du plan et vérifié l'égalité Débit/Crédit. Sans compte exact, la période est restée à vérifier et aucune ligne n'a été créée. \\
T08 & Demander à l'Assistant uniquement les documents bancaires. & La catégorie banque et les portées du cabinet et de l'entreprise sont appliquées. & Le moteur de requête a filtré explicitement la catégorie \texttt{banque}, appliqué \texttt{cabinet\_id} et \texttt{entreprise\_id}, et retourné des liens documentaires protégés. \\
T09 & Charger les routes fonctionnelles avec un administrateur de démonstration. & Chaque page authentifiée affiche un contenu réel sans erreur backend. & L'audit Chrome local a contrôlé 25 routes, notamment le nouvel accueil et la sélection d'entreprise ; le passage final est conforme. \\
T10 & Tenter des mutations après verrouillage d'une période, puis la réouvrir avec un rôle autorisé. & Les mutations sont refusées sans changement partiel ; les GET restent disponibles et la réouverture justifiée restaure les droits de modification. & Les routes de rapprochement bancaire, TVA, documents, reconstruction et clôture retournent HTTP~409 avant toute mutation sur une période verrouillée. Les consultations restent accessibles. Après réouverture autorisée et auditée, les mêmes opérations redeviennent possibles. \\
\bottomrule
\end{longtable}
\normalsize

\section{Contrôles de sécurité complémentaires}

Des contrôles ciblés ont été exécutés sur le code réel sans modifier
l'application ni les données de démonstration. Trois sondes ont sollicité
directement les dépendances de sécurité et le validateur d'import. En
complément, huit tests Pytest existants ont été exécutés avec succès : six cas
de sécurité documentaire, un cas d'isolation multi-tenant et un cas de
sélection des erreurs ouvrant droit à un retry.

\begin{table}[H]
\centering
\caption{Résultats des contrôles de sécurité complémentaires}
\footnotesize
\TableSetup
\begin{tabularx}{\textwidth}{|L{1.2cm}|L{4.3cm}|L{3.8cm}|Y|}
\toprule
\rowcolor{ComptaBlue}\TableHead{ID} & \TableHead{Méthode} & \TableHead{Résultat attendu} & \TableHead{Résultat réel / preuve} \\
\midrule
S01 & Modifier la signature d'un JWT valide, puis appeler la dépendance d'accès. & Rejet de l'authentification. & \TableSuccess{Conforme} : HTTP~401, jeton invalide ou expiré. \\
S02 & Présenter le rôle Assistant à la dépendance réservée à l'Expert-comptable ou au Chef de mission. & Refus de l'opération sensible. & \TableSuccess{Conforme} : HTTP~403, rôle insuffisant. \\
S03 & Exécuter le test de portée métier avec un autre \texttt{cabinet\_id}. & Refus d'un objet hors cabinet. & \TableSuccess{Conforme} : test Pytest réussi ; T03 conserve en outre la preuve HTTP~404 au niveau documentaire. \\
S04 & Soumettre un contenu PDF annoncé comme JPEG avec l'extension \texttt{.jpg}. & Refus avant stockage. & \TableSuccess{Conforme} : HTTP~400 ; les tests de taille, signature et confinement du chemin ont également réussi. \\
\bottomrule
\end{tabularx}
\end{table}

Ces contrôles sont inclus dans la suite complète de 401 tests backend validée ;
ils complètent les scénarios métier sans remplacer les tests d'intégration.

\section{Traçabilité des besoins fonctionnels}

La matrice suivante relie chaque besoin fonctionnel aux principaux éléments de réalisation ou de validation observables dans la version documentée.

\small
\TableSetup
\begin{longtable}{|L{1.6cm}|L{3cm}|L{9.8cm}|}
\caption{Matrice de traçabilité des besoins fonctionnels}\label{tab:tracabilite-bf}\\
\toprule
\rowcolor{ComptaBlue}\TableHead{Besoin} & \TableHead{État} & \TableHead{Preuve / validation} \\
\midrule
\endfirsthead
\toprule
\rowcolor{ComptaBlue}\TableHead{Besoin} & \TableHead{État} & \TableHead{Preuve / validation} \\
\midrule
\endhead
\TableId{BF01} & \TableSuccess{Réalisé} & Authentification, rôles et contrôles d'accès ; sondes S01 et S02 : HTTP~401 et 403 attendus. \\
\TableId{BF02} & \TableSuccess{Réalisé} & Isolation par cabinet et entreprise ; T03 : ressource d'un autre cabinet masquée par HTTP~404 ; contrôle S03 réussi. \\
\TableId{BF03} & \TableSuccess{Réalisé} & Import documentaire ; S04 et tests Pytest : format, taille, signature, extension, MIME et chemin contrôlés. \\
\TableId{BF04} & \TableSuccess{Réalisé} & Pipeline Groq/OCR et mécanismes de repli. \\
\TableId{BF05} & \TableSuccess{Réalisé} & Identification entreprise/tiers et classification achat/vente. \\
\TableId{BF06} & \TableSuccess{Réalisé} & Chronos numériques et filtres par période et catégorie. \\
\TableId{BF07} & \TableSuccess{Réalisé} & T01 : pré-écriture équilibrée ; T02 : incohérence signalée et aucune ligne déséquilibrée produite. \\
\TableId{BF08} & \TableSuccess{Réalisé} & Plan comptable propre à chaque entreprise et résolution des comptes. \\
\TableId{BF09} & \TableSuccess{Réalisé} & T04 : 2\,400 MAD alloués sur 4\,800 MAD et solde de 2\,400 MAD conservé. \\
\TableId{BF10} & \TableSuccess{Réalisé} & T05 : cours historique introuvable signalé, sans substitution par le cours courant. \\
\TableId{BF11} & \TableSuccess{Réalisé} & Grand Livre, Balance, TVA, CPC et Bilan. \\
\TableId{BF12} & \TableSuccess{Réalisé} & T06 : écart de bilan de 20 MAD signalé sans ligne fictive d'équilibrage. \\
\TableId{BF13} & \TableSuccess{Réalisé} & Tableau de bord, tâches planifiées avec heure et récurrence, messagerie, notifications, Assistant sécurisé et journal d'audit administrateur. \\
\TableId{BF14} & \textcolor{ComptaSlate}{\textbf{Partiellement réalisé}} & Exports génériques CSV, XLSX et PDF disponibles ; connecteur et intégration automatique directe avec Topaze non réalisés. \\
\bottomrule
\end{longtable}
\normalsize

\section{Problèmes rencontrés et corrections}

\subsection{Contraintes de ressources pour les traitements OCR et IA}

\textbf{Problème constaté.} L'une des principales difficultés rencontrées au
cours du développement concernait les ressources matérielles nécessaires aux
traitements documentaires exécutés localement. L'utilisation de PaddleOCR et,
surtout, l'exécution des modèles d'intelligence artificielle avec Ollama
pouvaient entraîner une consommation importante de mémoire vive et ralentir le
traitement sur ma machine de développement.

\textbf{Cause.} Ces traitements nécessitent le chargement local de modèles et
mobilisent des ressources pendant leur exécution. Cette contrainte était
d'autant plus importante que je ne disposais ni d'un serveur dédié ni d'une
infrastructure cloud personnelle suffisamment puissante pour héberger et
exécuter continuellement des modèles lourds.

\textbf{Solution adoptée.} J'ai intégré Groq comme service distant dans le
pipeline documentaire. Groq Vision constitue le chemin prioritaire pour la
classification et l'extraction structurée à partir des pages prétraitées. S'il
n'est pas disponible ou si son résultat n'est pas exploitable, le système tente
d'abord l'extraction native du texte d'un PDF, puis utilise le prétraitement
OpenCV et PaddleOCR lorsque l'OCR est nécessaire. Le texte obtenu peut ensuite
être analysé par Groq texte, avec une branche spécialisée pour les relevés
bancaires.

\textbf{Résultat et continuité du traitement.} Cette organisation limite le
recours systématique aux modèles exécutés sur la machine locale, sans supprimer
les mécanismes de repli. En cas d'indisponibilité des services Groq, le pipeline
conserve l'extraction locale par PaddleOCR ainsi que les expressions régulières,
les heuristiques de classification et l'enrichissement sémantique partiel par
Ollama. La solution réduit ainsi qualitativement la charge locale tout en
préservant la flexibilité du pipeline et la continuité du traitement, sans
supposer qu'un résultat ambigu puisse être validé automatiquement.

% ============================================================
% 13. Conclusion
% ============================================================
\chapter{Conclusion générale}

\section{Bilan}

ComptaFlow fournit une chaîne cohérente allant du document brut aux états comptables. Son principal apport réside dans l'association de techniques d'extraction variées avec des règles métier déterministes et un contrôle humain. La structure multi-tenant, la conservation des données originales, la résolution des comptes par entreprise et le refus des corrections fictives répondent aux exigences de confiance du domaine comptable.

Le périmètre réalisé couvre l'import, les chronos, la préparation et la validation des pré-écritures, le workflow pré-Topaze, la gestion bancaire avancée, les devises BAM, les états comptables, la TVA, les contrôles et le verrouillage de pré-clôture. Il comprend également l'Assistant sécurisé, le journal d'audit, la messagerie et la planification des tâches. La modularité des services permet d'ajouter des traitements tout en maintenant des règles identiques entre le chemin principal et les replis.

Les résultats obtenus répondent ainsi à la problématique initiale en réduisant les opérations manuelles situées en amont de la saisie comptable tout en préservant la fiabilité et le contrôle professionnel. La version documentée a validé 401 tests backend, 18 tests frontend, le typecheck, le lint et le build de production, ainsi qu'un audit visuel de 25 routes authentifiées. Cette automatisation reste encadrée par une validation humaine pour les décisions comptables sensibles. L'intégration automatique avec Topaze, l'évaluation quantitative des performances OCR/IA et le renforcement de l'industrialisation, des sauvegardes et de l'observabilité constituent les principales étapes à poursuivre avant une utilisation en production.

\section{Compétences acquises}

La réalisation de ce projet m'a permis de renforcer mes compétences en analyse
des besoins et en compréhension d'un processus métier réel. L'étude du
fonctionnement comptable de SEGURIBAT m'a amenée à traduire un processus encore
largement manuel en besoins fonctionnels, règles métier et modèles de données
exploitables dans une application. J'ai ainsi appris à relier les observations
du terrain aux décisions de conception, puis à faire évoluer ces décisions à
partir des retours obtenus durant le stage.

Sur le plan du développement, j'ai approfondi mes compétences en conception et
réalisation d'applications web avec React, TypeScript et Vite pour le frontend,
ainsi qu'avec FastAPI, SQLAlchemy et PostgreSQL pour le backend et la
persistance des données. Ce travail m'a permis de mieux maîtriser la
structuration d'une architecture modulaire, les échanges avec une API REST,
l'authentification, la gestion des rôles et le suivi des évolutions du schéma de
données avec Alembic.

Le projet m'a également permis d'aborder des problématiques spécialisées de
traitement documentaire. J'ai appris à intégrer et à articuler plusieurs
mécanismes d'extraction et d'analyse, notamment Groq, PaddleOCR et Ollama, tout
en prévoyant des méthodes de repli adaptées aux limites matérielles et aux
différences de qualité des documents. L'utilisation de Celery et Redis m'a
également permis de mieux comprendre l'exécution asynchrone des traitements
longs et la séparation entre la réception d'une requête et son traitement en
arrière-plan.

Enfin, j'ai acquis une meilleure compréhension des contraintes propres au
domaine comptable : conservation des données originales, préparation des
pré-écritures, résolution des comptes à partir du plan de chaque entreprise,
rapprochement bancaire, contrôles de cohérence et maintien d'une validation
humaine pour les décisions sensibles. L'exécution des tests backend et des
contrôles frontend, l'analyse des anomalies et les corrections successives ont
renforcé ma capacité à valider une application de manière structurée et à
justifier mes choix techniques et fonctionnels.

\section{Perspectives}

ComptaFlow constitue une base fonctionnelle qui pourrait être prolongée par des
évolutions ciblées, sans remettre en cause le contrôle humain placé au cœur du
traitement comptable :

\begin{itemize}[leftmargin=*]
  \item Une évolution prioritaire consisterait à développer un mécanisme
  d'export dans un format compatible avec Topaze ou, si des possibilités
  techniques sont confirmées ultérieurement, une intégration avec ce logiciel.
  ComptaFlow continuerait ainsi à automatiser la réception, le classement,
  l'extraction, le contrôle et la préparation comptable, tandis que Topaze
  resterait l'outil comptable de référence utilisé par SEGURIBAT. Aucune
  intégration automatique directe avec Topaze n'a été réalisée durant le stage.

  \item Un corpus de documents anonymisés et versionnés pourrait être constitué
  afin de mesurer la qualité des différentes méthodes d'extraction et de
  classification. Cette évaluation permettrait de comparer les résultats selon
  le type et la qualité des pièces, sans présumer d'un taux de précision qui n'a
  pas encore été mesuré.

  \item Les mécanismes d'assistance pourraient enfin être enrichis afin de mieux
  classer les situations ambiguës et de présenter des choix plus pertinents à
  l'utilisateur, tout en conservant une validation humaine pour les décisions
  comptables sensibles.
\end{itemize}

\subsection{Conditions préalables à une mise en production}

La version documentée constitue une base fonctionnelle, mais son exploitation
continue nécessiterait une phase d'industrialisation. Les conditions
prioritaires sont les suivantes :

\begin{itemize}[leftmargin=*]
  \item organiser les sauvegardes de PostgreSQL et des documents, puis tester
  régulièrement une restauration complète ;
  \item déployer les échanges sous HTTPS/TLS et mettre en place une gestion et
  une rotation maîtrisées des secrets ;
  \item superviser l'API, les workers Celery, Redis et PostgreSQL, ainsi que
  l'espace disque occupé par le stockage documentaire ;
  \item suivre les temps de traitement, les erreurs d'extraction et les
  mécanismes de repli utilisés afin de détecter les dégradations ;
  \item définir une politique de conservation et de suppression, ainsi que les
  règles de confidentialité applicables aux données transmises aux services
  externes ;
  \item compléter la validation par des tests end-to-end et constituer un corpus
  OCR/IA anonymisé et versionné pour les évaluations quantitatives.
\end{itemize}

% ============================================================
% 14. Références
% ============================================================
\begin{thebibliography}{99}
\addcontentsline{toc}{chapter}{Références}

\bibitem{fastapi}
FastAPI, \textit{Documentation officielle}, \url{https://fastapi.tiangolo.com/}.

\bibitem{sqlalchemy}
SQLAlchemy, \textit{Documentation officielle}, \url{https://docs.sqlalchemy.org/}.

\bibitem{alembic}
Alembic, \textit{Documentation officielle}, \url{https://alembic.sqlalchemy.org/}.

\bibitem{react}
React, \textit{Documentation officielle}, \url{https://react.dev/}.

\bibitem{typescript}
Microsoft, \textit{Documentation TypeScript}, \url{https://www.typescriptlang.org/docs/}.

\bibitem{postgresql}
PostgreSQL Global Development Group, \textit{PostgreSQL Documentation}, \url{https://www.postgresql.org/docs/}.

\bibitem{celery}
Celery Project, \textit{Celery Documentation}, \url{https://docs.celeryq.dev/}.

\bibitem{redis}
Redis, \textit{Redis Documentation}, \url{https://redis.io/docs/}.

\bibitem{paddleocr}
PaddlePaddle, \textit{PaddleOCR Documentation}, \url{https://www.paddleocr.ai/}.

\bibitem{groq}
Groq, \textit{Groq API Reference}, \url{https://console.groq.com/docs/api-reference}.

\bibitem{ollama}
Ollama, \textit{Documentation officielle}, \url{https://docs.ollama.com/}.

\bibitem{docker}
Docker, \textit{Docker Documentation}, \url{https://docs.docker.com/}.

\bibitem{bam}
Bank Al-Maghrib, \textit{Cours de change}, \url{https://www.bkam.ma/}.

\bibitem{cgnc}
Conseil National de la Comptabilité, \textit{Code Général de Normalisation Comptable marocain}.

\bibitem{repo}
Projet ComptaFlow, \textit{Code source, modèles, services, routes, migrations et tests de la version documentée}, dépôt local consulté en août 2026.

\end{thebibliography}

% ============================================================
% 15. Annexes
% ============================================================
\appendix
\renewcommand{\chaptername}{Annexe}
\chapter{Arborescence technique}

\begin{lstlisting}
ComptaFlow_vs/
|-- backend/
|   |-- alembic/versions/       migrations de base
|   |-- app/
|   |   |-- api/                routes FastAPI
|   |   |-- core/               configuration, securite, DB, Celery
|   |   |-- models/             entites SQLAlchemy
|   |   |-- schemas/            contrats Pydantic
|   |   |-- services/           logique metier
|   |   `-- tasks/              traitements asynchrones
|   `-- tests/                  tests Pytest
|-- frontend/
|   `-- src/
|       |-- api/                clients HTTP
|       |-- components/         composants reutilisables
|       |-- context/            authentification
|       |-- pages/              ecrans fonctionnels
|       |-- types/              types TypeScript
|       `-- utils/              utilitaires
|-- docker-compose.yml
`-- rapport_comptaflow.tex
\end{lstlisting}

\chapter{Commandes de validation}

\begin{lstlisting}
# Backend (depuis backend/)
python -m pytest -q
python -m alembic heads
python -m alembic current
python -m alembic history

# Frontend (depuis frontend/)
npm run typecheck
npm run build
npm run lint

# Compilation du rapport (depuis la racine)
pdflatex rapport_comptaflow.tex
pdflatex rapport_comptaflow.tex

# Variante si latexmk est disponible
latexmk -pdf rapport_comptaflow.tex
\end{lstlisting}

\end{document}
